\documentclass[../main.tex]{subfiles}
\graphicspath{{\subfix{../}}}
\begin{document}

\textit{Fotoeffektus, Compton-effektus. a hőmérsékleti sugárzás spektruma (beleértve: Stefan–Boltzmann-törvény, Wien-törvény), Rutherford-kísérlet, Millikan-kísérlet,
Davisson–Germer-kísérlet, Stern–Gerlach-kísérlet, Einstein–de Haas-kísérlet, Zeeman-effektus.}

\subsection{Fotoeffektus}
H. Hertz és munkatársai a szikrakisülés tanulmányozásakor észrevették, hogy a kisülés könnyebben létrejön, ha fénnyel világítják meg a katódot.
Ez adta a lökést a fény kvantumos szerkezetére vonatkozó egyik legfontosabb kísérleti tény felfedezéséhez. Ez a H. Hertz által 1887-ben felfedezett fényelektromos hatás, más néven fotoelektromos hatás vagy fotoeffektus.
A fényelektromos hatást P. Lenard, majd mások is tanulmányozták.
\begin{figure}[H]
    \centering
    \includegraphics[width=0.5\textwidth]{imgs/10-tetel/fotoeffektus.png}
    \caption{Fotoeffektus kísérleti elrendezése}
\end{figure}
Egy vákuumba elhelyezett kondenzátor egyik fémfegyverzetét látható fénnyel megvilágítjuk, és a körben folyó áramot érzékeny galvanométerrel mérjük. Ezt az elrendezést - mai felhasználása alapján - nevezhetjük fotocellának is. A kísérlet eredménye, hogy
a fény hatására a fémlapból elektronok lépnek ki, és ezek I áramot okoznak a galvanométeren. A kondenzátorra kapcsolt V feszültséggel lehet az elektronokat gyorsítani, ekkor a $V$-tól független áramot mérünk, adott megvilágításnál.
Ha azonban lassító feszültséget kapcsolunk a kondenzátorra, akkor egy $V_0$ ellenfeszültségnél megszűnik az áram, ezt nevezzük lezáró feszültségnek.

Ha nagyobb intenzitású fényt alkalmazunk, akkor a lezáró feszültség változatlan, de a kilépő elektronok intenzitása megnő, így az $I$ áramerősség $I_2 > I_1$. A lezáró feszültség magyarázata az, hogy
$V$ lassító feszültség $e \cdot V$ energiával lassítja le az elektront, ha a kezdeti mozgási energiája ennél kisebb volt, akkor mielőtt a túlsó kondenzátor lemezre érkezik, a sebessége O-ra csökken, az elektron visszafordul, nem jut át, így
ezen elektronok nem vesznek részt az áramvezetésben. Ha a V éppen akkora, hogy a fémlapról induló maximális sebességű elektronokat is visszafordítja, akkor az áramvezetés itt megszűnik: $\frac{1}{2}mv_{\text{max}}^2 = eV_0$. Tehát a lezáró feszültség arányos a kilépő elektronok
mozgási energiájával, így azok sebességnégyzetéveL Ha megváltoztatjuk a fémlemez anyagát vagy a megvilágító fény színét (frekvenciáját) , akkor a lezáró feszültség is változik, viszont (ahogy láttuk) a fény intenzitásától a
lezáró feszültség független.
\begin{figure}[H]
    \centering
    \includegraphics[width=0.5\textwidth]{imgs/10-tetel/fotoeffektus2.png}
    \caption{Fotoeffektus lezáró feszültség és frekvencia kapcsolata}
\end{figure}

Ha egy kísérletsorozatot hajtunk végre adott fémlappal, de különböző frekvenciájú fénnyel, és mérjük a lezáró feszültséget, akkor lineáris összefüggést kapjuk. (Ezt a kísérletet először Millikan végezte
el 1910-ben!) Ábrázoljuk az $x$ tengelyen a fény frekvenciáját $v$-t , és az $y$ tengelyen az ezekhez a frekvenciákhoz kapott lezáró feszültséget az elemi töltéssel megszorozva 
(így $y = eV_0 = \frac{1}{2}mv_{\text{max}}^2$)· A kísérleti tapasztalat az, hogy egy egyenest alkotnak a mérési pontok, az egyenes egy adott $v_0$ küszöbfrekvenciánál metszi az $x$ tengelyt, 
ennél kisebb frekvenciájú fényesetén nem kelllezáró potenciált alkalmazni, a galvanométeren ilyenkor egyáltalán nem mérl1ető áram. A másik érdekes tapasztalat, hogy az egyenes meredeksége megegyezik a fekete test sugárzásakor definiált $h$ Planck-állandóval.

Tehát a tapsztalatokat összefoglalva:
\begin{enumerate}
    \item A fénnyel megvilágított fémből a fény hatására elektronok lépnek ki.
    \item A kilépő elektronok száma arányos a beeső fény intenzitásával.
    \item A kilépő elektronok sebessége a fény frekvenciájától függ , és nem függ a beeső fény intenzitásától.
    \item Adott frekvenciájú megvilágítás esetén a kilépő elektronok sebessége folytonosan változhat egy maximális sebességig, amely nagysága a megvilágító fény frekvenciájától és a fém anyagától függ.
    \item Az elektronok áramának elindításához egy küszöbértéknél nagyobb frekvenciájú fény szükséges.
    \item Az elektronok árama a megvilágítás után azonnal megindul.
\end{enumerate}

A fényelektromos hatás fenti törvényszerűségeit nem lehetett a klasszikus hullámelmélettel magyarázni. A hullámelmélet alapján ugyanis a következő magyarázatot lehetne elképzelni: férore beeső elektromágneses hullám a
fém elektronjait rezgésbe hozza. Ha a beeső hu ll ám frekvenciája megegyezik az elektron sajátrezgésének frekvenciájával, az elektron rezgési amplitúdója olyan nagy lesz, hogy az kiszakad a fémből. A fémből kilépő elektronok a
mozgási energiájukat a beeső fényhullámból veszik, ezért azt várjuk, hogy a fotoelektronok energiája a fény intenzitásától függjön. (Nagyobb intenzitású fényhullámban nagyobb az elektromos térerősség $I \sim E^2$, így nagyobb
erő hozza rezgésbe az elektronokat.) A kísérletek eredményei szerint azonban a kilépő elektronok energiája teljesen független a beeső fény intenzitásától, és csak a fotoelektronok száma arányos ezzel. A fotoelektronok sebessége
(energiája) a beeső fény frekvenciájától függ, mégpedig nő a frekvenciával.

Felmerült az a gondolat is, hogy az elektronok a kilépéshez szükséges energiát a fényhullám energiájából fokozatosan gyűjtik össze. Ennek az elképzelésnek a helytelenségét azonnal be lehet látni: a beeső hullám energiája
megoszlik a fém nagyszámú elektronja között , így kis fényintenzitásnál tekintélyes idő kellene ahhoz, hogy az elektronok összegyűjtsék a kilépéshez szükséges energiát. A fémfelület által elnyelt energia a fényerősséggel,
a megvilágított felület nagyságával és a besugárzási idővel arányos. A fényelektromos hatás igen kis felület megvilágításánál is létrejön: $10^{-8} W/m^2$-es felületi megvilágítás esetén-ezt hozza létre 100 km távolságból egy 100 W-
os égő - száz óra alatt nyel ődik el annyi fény, amennyi egy fényelektron kilépéséhez szükséges, azaz néhány $eV$. Ezzel szemben a kísérletek azt mutatják, hogy az elektronemisszió a megvilágítást követően $10^{-9} s$-on belül
létrejön. (A holdfény intenzitása $\sim 10^{-3} W/m^2$. Ha $1 \text{\AA}$ rácsállandójú fémmel számolunk, akkor $1 m^2$-en $10^{20}$ atom foglal helyet a fém felületén, ha mindegyik egy vegyértékelektronja kapja a be es ő fény intenzitásá nak ráeső részét , akkor 104 s, azaz néhány óra alatt kap egy elektron $1 eV$, azaz
$10^{-19} J$ energiát. De a tapasztalat szerint a holdfény is képes fotoeffektust létrehozni.) A fotoeffektus kísér le ti ténye tú lmutat a klasszikus fizikán, és magyarázatát 1905-ben találták meg, amikor a modern fizika szemlélete -
Planck hipotézi sé ből kiindulva - kezdett elfogadottá válni.

\subsubsection{A fotoeffektus kvantumos magyarázata: a fotonkép}
A fényelektromos hatás magyarázatát A. Ei nste in adta meg 1905-ben, Planck kvantumhipotézise általánosításávaL Rámutatott, hogy a nehézségek mind megszűnnek , ha korpuszkuláris álláspontra helyezkedünk, és a fényt
$h \nu$ energiájú kvantumok, ún. fotonok áramának fogjuk fel. Az elnyelt foton átadja energiáját az elektronnak, és ha az energia e legendő ahhoz, hogy kiszabadítsa az elektront a visszatartó erők kötelékéből , akkor az elektron
kilép a fémből. Annak valószínűsége , hogy egy elektron két fotont nyeljen el, igen kicsi, minden kilépő elektron egy foton energiáját viszi el. (Később ézerrel előállított intenzív fénysugárral sikerült két fotonos fotoeffektust
létrehozni.) Az így kivál tott elektronok számának az elnyelt foton ok számával, vagyis a fény intenzitásával kell arányosndk lenni, és a kísérletek is ezt mutatták. Ebben a képben a korábban említett kísérletsorozat, a különböző
frekvenciájú fényekkel. A tapasztalati $eV_0 = h(\nu - \nu_0)$ összefüggést kicsit átírjuk , ahol $V_0$ a lezáró potenciál, és $\nu_0$ a küszöbhullámhossz.
Felhasználva, hogy $eV_0 = \frac{1}{2} mv_{\text{max}}^2$ , az egyenletet a következő alakban írhatjuk fel:
\begin{equation}
    h \nu = \frac{1}{2} mv_{\text{max}}^2 + h \nu_0
\end{equation}
Itt bevezethetjük a munkafüggvény fogalmát: $W = h \nu_0$. Ez az a minimális energia, amely ahhoz szükséges, hogy az elektron a fémből kiszabaduljon. Ezzel az egyenlet a következő alakot ölti:
\begin{equation}
    h \nu = \frac{1}{2} mv_{\text{max}}^2 + W
\end{equation}
A fotonképben egy foton egy elektront üt ki a fémből, és a folyamat során az energia megmaradását éppen ez az egyenlet írja le (Einstein-egyenlet). Ez azt jelenti, hogy egy foton energiája:
\begin{equation}
    E_{\text{foton}} = h \nu
\end{equation}
Ebből az is látszik, hogy a fémből egy elektron kiszakításához legalább $W$ munka szükséges. Ez a kilépési munka az egyes fémekre más és más érték, de
álta lában $eV$ nagyságrendű . Az elektronok nemcsak a maximális sebességgel lépnek ki a fémből, hanem sebességeloszlásuk folytonos , ami azt mutatja,
hogy vannak kevésbé, ill. mélyebben kötött elektronok a szilárd testek fémes rácsában. Ezt a mai szóhasználattal elektrontengernek nevezzük, és $W$
jelentése, hogy az elektrontenger felszínén $W$ energiával kötött elektronok helyezkednek el. (Az elektrontenger igazi jelentését a modern szilárdtestfizika adja meg a Pauli-elv segítségével)
Megjegyzendő, hogy a fotoeffektus csak az egyik módja elektronok kiütésének egy fém lemezről, más folyamatok is ki tudják szabadítani a fémes rácsban kötött elektronokat. Ezek közül hármat sorolunk fel:
\begin{itemize}
    \item Termikus elektronemisszió, amikor addig melegítik a fémet, míg felületéről az elektronok a hő mozgás miatt ki tudnak lépni. Ilyenkor az elektronok $\sim kT$ termikus energiája elegendő a $W$ kilépési munka legyőzésére.
    \item Hideg elektronemisszió, ilyenkor külső elektromos tér segít az elektronoknak, hogy kiszabaduljanak. Kisülési csövekben is ilyen folyamat zajlik le, és a kiszabaduló elektronok ionizálva a gázt létrehozzák a kisülést . Katódsugárcsőben a nagy elektromos tér hatására szabadulnak ki az elektronok a fémből a vákuumba.
    \item Másodlagos sugárzás, ilyenkor más részecskékkel bombázzák a fém felületét, és ezek adják át mozgási energiájukat az elektronnak, amely így ki tud szabadulni a fémből.
\end{itemize}

\subsubsection{Atomi fotoeffektus}
Történetileg a fotoeffektust először fémlemezen figyelték meg, de fizikailag ugyanaz a folyamat végbemegy különálló atomokon is. Ha fénnyel vagy még
rövidebb hullámhosszú fotonokkal bombázunk atomokat, a fotonok kitudnak lökni egy-egy elektront az atomból, egyben a bombázó foton megsemmisül. Ilyenkor teljesen azonos egyenleteket tudunk alkalmazni, csak a $W$
kilépési munka helyett az elektron kötési energiáját kell fogalmilag használni. Az atomi fotoeffektus nagyon fontos folyamat az elektromágneses hullámok és az anyag kölcsönhatásának leírásánál, de hozzá az atom szerkezetének
részletesebb ismerete szükséges, melyet a következő fejezetekben tárgyalunk csak. Megjegyezzük, hogy szabad elektron nem tud elnyelni fotont, azaz fotoeffektus nem megy végbe szabad elektronon. Ennek ellenkezője áll fenn:
minél kötöttebb egy atomi elektron, annál nagyobb a valószínűsége, hogy egy (megfelelően) adott energiájú foton fotoeffektussal ki tudja ütni az atom kötéséből.

\subsection{Compton-effektus}
A fotoeffektus so rán az elektromágneses hullámok (például fénysugarak) elektront ütnek ki a fémből vagy különálló atomokból, de ilyenkor a fotonmegsemmisül (elnyelődik). A foton-elektron kölcsönhatásnak van
olyan formája is, amikor a beeső foton kölcsönhat az elektronnal, de nem semmisül meg, ez a Compton-effektus.

\subsubsection{Compton kísérletei a röntgensugárzás rugalmatlan szóródására}
1923-ban A. Campton mutatta ki, hogy van olyan folyamat is, amikor fotonok elektronokon szóródnak, és a szóródott fotonak frekvenciája nem egyezik
meg a beeső fotonak frekvenciájával. Itt ismét egy olyan folyamatra bukkantunk, ami nem magyarázható meg tisztán a klasszikus elektrodinamika
segítségéveL Az ugyanis azt mondja, hogy egy elektromágneses hullám, ha ráesik egy elektronra , akkor rezgésbe hozza azt (rugalmasan kötött elektronmodell).
A rezgő elektron a gyorsulása miatt sugároz, de a rezgésének frekvenciája megegyezik a gerjesztő rezgés frekvenciájával, így a kisugárzott
hullámok is azonos hullámhosszúak lesznek. Ez a fotonképben annyit jelent, hogy a beeső foton és a szóródott foton frekvenciája, így energiája is, azonos. Ez a Thomson-szórás.
A klasszikus elektrodinamika csak rugalmas szóródást jósol.
\begin{figure}[H]
    \centering
    \includegraphics[width=0.5\textwidth]{imgs/10-tetel/compton.png}
    \caption{Compton-effektus kísérleti elrendezése}
\end{figure}
Compton kísérletében egy molibdén antikatódú ($A$) röntgencsővel ($R$) adott frekvenciájú röntgensugárzást állított elő (lásd később a karakterisztikus röntgensugárzásnál), és azt elektronban gazdag szóróközegre ($Sz$), grafitra irányította.
Az adott szögben kollimátorokat ($K$) helyezett el, melyeken az átmenő sugárzás hu ll ámhosszát egy visszaverő kristállyal ($Kr$) és
egy ionizációs kamra ($IK$) segítségével mérte. A hullámhossz mérése úgy történik, hogy a koBimátorokon átjövő fotonok ráesnek a kristályra ($CaCO_3$)
és az interferenciának megfelelő szögben visszaverődnek a kristályról. A mozgatható ionizációs kamra a visszaverődés szögében fog maximális áramot jelezni, amiből a $2d\sin \phi = n\lambda$ 
összefüggés alapján a hullámhossz kiszámolható. A röntgencső forgatásával lehet változtatui a $\vartheta$ szórási szöget. Compton kísérleteinek eredményei:
\begin{enumerate}
    \item A rugalmatlan szóródás felfedezése, ami mindig a hullámhossz növekedésével jár.
    \item A hullámhossz-növekedés az eltérülési szöggel növekszik; Compton kimérte a $\Delta \lambda$ változását $\vartheta$ szög szerint.
    \item A hullámhossz növekedése adott szögben független a szóróközeg anyagától. Ez a tapasztalat mutatta meg, hogy nem az atomok, hanem csak azok individuálisnak tekinthető elektronjai vesznek részt a folyamatban.
\end{enumerate}
Ezt a kísérletet azért kellett röntgensugarakkal elvégezni, mert a rugalmatlan szórás csak szabad- vagy kváziszabad-elektronok esetében lép fel,
amikor a foton energiája sokkal nagyobb az elektronok kötési energiájánál ($h \nu \gg E_k$)·

\subsubsection{A Compton-effektus értelmezése}
Ezeket a kísérleti tapasztalatokat könnyen lehet magyarázni, ha feltételezzük, hogy a sugárzás korpuszkuláris szerkezetű, azaz diszkrét energiacsomagok - fotonok - árama. Tegyük fel, hogy egy $h \nu$ energiájú, $p = \frac{h \nu}{c}$ lendületű
foton esik egy elektronra. Az ütközés után a foton $\vartheta$ szögben repül ki $\lambda$ megváltozott hullámhosszal, ill. $h \nu'$ lecsökkent energiával. Közben az elektron $\phi$ szögben visszalökődik, lendületre és mozgási energiára tesz szert. Ha
felírjuk az energia és az impulzus megmaradását, a hullámhosszváltozás kiszámolható.
\begin{figure}[H]
    \centering
    \includegraphics[width=0.5\textwidth]{imgs/10-tetel/compton2.png}
    \caption{Compton-effektus ütközési sémája}
\end{figure}
Az energiamegmaradás relativisztikusan írandó:
\begin{equation}
    h \nu + m_0 c^2 = h \nu' + \sqrt{p_e^2 c^2 + m_0^2 c^4}
\end{equation}
Itt $m_0 = 511 \, \text{keV}/c^2$ az elektron nyugalmi tömege, $p_e$ pedig az elektron lendületének nagysága, $\nu$ és $\nu'$ a foton beesési és kilépési frekvenciája.
A lend ület vektormennyiség, így két irányú megmaradását is felírhatjuk, most mégis célszerű a vektoregyenl etet felírni, és a lendületmegmaradást
a koszinusztétellel kifejezni. A $p_{be}$ és $p_{ki}$ a foton beeső és ütközés utáni lendületét jelöli:
\begin{equation}
    p_{be} = p_{be}^2 + p_{ki}^2 - 2 p_{be} p_{ki} \cos \vartheta
\end{equation}
Ezt $c^2$-el végígszorozva és a foton lendületére fennálló $\frac{h \nu}{c}$ formulát alkalmazva:
\begin{equation}
    p_e^2 c^2 = h \nu^2 + h \nu'^2 - 2 h^2 \nu \nu' \cos \vartheta
\end{equation} 
Ha az energiamegmaradás egyenletében a $h \nu'$-t átvisszük a másik oldalra, majd négyzetre emelünk és $p_e^2 c^2$ helyére a lendületmegmaradás egyenletéből adódó értéket írjuk,
akkor csak $\nu'$ marad ismerretlen, sőt a négyzetes tagok kiesnek, és a következő formulát kapjuk a szórt foton energiájára:
\begin{equation}
    h \nu' = \frac{h \nu}{1 + \frac{h \nu}{m_0 c^2} (1 - \cos \vartheta)}
\end{equation}
Ebből a $\Delta = \frac{c}{\nu'} - \frac{c}{\nu}$ alapján:
\begin{equation}
    \Delta \lambda = \lambda' - \lambda = \frac{h}{m_0 c} (1 - \cos \vartheta)
\end{equation}
adódik. Itt a $\frac{h}{m_0 c}$ univerzális állandók hosszúság-dimenziójú kombinációja, neve az elektron Compton-hullámhossza (értéke 2,426 pm). 
Ilyen hullámhosszúságú foton energiája éppen az elektron nyugalmi energiájával egyezik meg. Érdemes megjegyezni, hogy a hullámhossz-eltolódás a Compton-
hullámhosszon kívül csak a szögtő l függ, a beeső foton hullámhosszától nem, ahogy ezt Compton kísérleteiben is tapasztalta.
A formulák pontosan leírják a Compton-kísérletben kapott eredményeket, így alátámasztják a foton szabadelektronokon történő szóródásának ezen részecskemodelljét.
A Compton-effektus egyben megmutatja, hogy az energia- és a lendületmegmaradás elemi foton-elektron folyamatokban is megmarad, nem csak valamilyen nagy statisztikus sokaság átlaga esetén. Megjegyezzük,
hogy ez a gondolatmenet nem zárja ki a fotonok hullámszerű viselkedését a Compton-effektus során, de a részecskekép sokkal elegánsabban magyarázza
ezt az effektust. 1927-ben E. Schrödinger a klasszikus elektrodinamika és a későbbiekben tárgyalandó anyaghullámok segítségével adott egy - a jelen munkát meghaladó - hullámleírást a fotonok rugalmatlan szóródására.
A Compton-effektus tulajdonságai:
\begin{itemize}
    \item Az elektronnak átadott energia lehet nagyon kicsiny, mikor a foton alig, térül el, és van egy maximális értéke, amikor a foton $\vartheta = 180\degree$-ban éppen visszafelé szóródik. Ilyenkor $h \nu - h \nu' = \frac{h \nu}{\frac{m_0 c^2}{2 h \nu} + 1}$ energiát visz el az elektron, s mindig marad valamennyi a fotonnak.
    \item Kis frekvenciájú határesetben a Campton-effektus átmegy a klasszikusan jól ismert Thomson -szórásba. Az összefüggései $h \nu \ll m_0 c^2$ esetén visszaadják a Thomsan-szórás klasszikusan levezetett formuláit.
\end{itemize}
Compton-effektus nehezebb elemi részecskéken is végbemegy, így például protonon, vagy egy teljes atomon. A proton tömege $\sim 2000$-szer nagyobb az elektron tömegénél, ezért a Compton-hullámhossza $\sim 2000$-szer kisebb, így a fm – azaz $10^{-15}\,\mathrm{m}$ – nagyságrendbe esik. Protonon a röntgensugarak hullámhosszváltozása elhanyagolható, hiszen a röntgensugarak hullámhossza pm-es nagyságrendű.

Compton eredeti kísérletében minden szórási szögnél két hullámhossz esetén kapott maximumot az ionizációs kamra áramában. Az egyik az elektronon történő rugalmatlan szóródáshoz tartozik, amiről eddig szó volt, a másik pedig a besugárzó hullámhosszhoz. Ezek a rugalmas szóródáshoz tartozó események, amikor a fotonok az egész atomon szóródnak. Az atom tömege a proton tömegének néhányszorosa, így a hullámhossz-növekedés elhanyagolható.
Az atom keresztmetszete (az elektronéval összehasonlítva) kicsi, de a rugalmatlan szóródás valószínűsége lényegesen nagy az atomon, mint egészben történő szóródásé.

\textit{Összefoglalva:} a Compton-effektussal olyan új folyamatot ismertünk meg, amelyben a foton energiája részben egy elektron mozgási energiájává alakul át, és egy kisebb energiájú foton is létrejön. A kísérletek igazolták, hogy az energia- és a lendületmegmaradás az egyes elemi folyamatokban is fennáll.
A folyamatot a kísérleti eredményekkel összhangban jól lehet a kvantumhipotézis segítségével értelmezni. A Compton-effektus vizsgálata még inkább előtérbe helyezi az elektromágneses sugárzás részecsketermészetét. Már nemcsak annyiban kell részecskének tekinteni a fotont, hogy $h\nu$ nagyságú energiakvantumot képvisel, hanem mint ami $h/\lambda$ lendületet is hordoz, és ezzel az összefüggéssel a lendületmegmaradás tétele a kísérletekkel megegyező eredményt ad.

Térjünk vissza röviden a fotoeffektushoz. Ott említettük, hogy szabad elektronon a foton nem tud elnyelődni. Ezt a tényt a foton impulzusának ismeretében, már meg is tudjuk alapozni. Ilyenkor a lendület megmaradása 
$p_e = p_{be}$
feltételt jelent, így az energia és a lendület egyszerre nem tud megmaradni, mert az elektron nyugalmi tömege nem 0. Ez nincs ellentmondásban, hanem éppen azt jelenti, hogy szabad elektron a foton energiáját, vagyis a lendületétől volt szó, amint azt említettük, hogy vele, vagy ha közelében egy rácspont található, amely a kristályrács, amely a felesleges impulzust elviszi. Ez jelzi, hogy milyen fontos szerepe van a harmadik résztvevő testnek a fotoeffektus során.

\subsection{Hőmérsékleti sugárzás spektruma}
\subsubsection{Az abszolút fekete test}
A kvantumfogalom megalkotásával összefüggő kutatásokban jelentős szerepet játszott egy idealizált fogalom, az ún. „abszolút fekete test” fogalma. Mint ismeretes, ha egy testre fényt bocsátunk, egy része visszaverődik, másik része bejut a testbe. A bejutott fény részben (esetleg teljesen) elnyelődik, részben kilép a test túlsó felületén.
A testeket számunkra az teszi láthatóvá, hogy visszavernek fényt. Az olyan test, amelyik a ráeső fényt teljesen elnyelné, abszolút fekete, azaz láthatatlan lenne. A fekete test által elnyelt sugárzás energiát visz be a testbe, megnő a test belső energiája, következésképpen magasabb lesz a hőmérséklete. A hőmérséklet korlátlanul növekedne,
ha nem létezne olyan folyamat, amely a belső energiát csökkenti. Ennek viszont olyannak kell lennie, hogy vákuumban is érvényesülni tudjon. A test belső energiájának egyensúlya sugárzás kibocsátása útján valósul meg. Ezen egyensúlyi sugárzást nevezzük \textit{hőmérsékleti sugárzásnak}, vagy a fekete test sugárzásának. Ennek törvényeinek vizsgálata vezetett el először a kvantumfogalomhoz.

Az abszolút fekete test a természetben nem létezik, de hasznos és jogos idealizációnak bizonyult. Kísérletileg előállítható olyan rendszer, amelynek segítségével meg lehet vizsgálni ezen idealizáció – az abszolút fekete test – tulajdonságait.

Olyan anyag, amelynek felülete abszolút fekete volna, vagyis minden ráeső sugárzást teljesen elnyelne, nincs. Megvalósítható azonban a fekete test úgy, hogy egy a sugárzást át nem eresztő, kormozott falú, belül üreges zárt edény falán kis lyukat fúrunk, ahogy azt a 3.1. ábra szemlélteti.
A lyukon át behatoló sugárzás ugyanis az üreg belső felületén való sokszoros, többnyire diffúz visszaverődés következtében gyakorlatilag teljesen abszorbeálódik, mielőtt a lyukból kijutna. Ha pl. egyszeri visszaverődés esetén a sugár energiájának csak $90\%$-a absorbeálódik is, a tizedik visszaverődés után a beeső energia $10^{-10}$-szerese marad már csak meg. Valóban, az ilyen lyukat sokkal feketébbnek látjuk, mint az edény bármennyire is bekormozott külső falát.
\begin{figure}[H]
    \centering
    \includegraphics[width=0.5\textwidth]{imgs/10-tetel/fekete_test.png}
    \caption{Az abszolút fekete test modellje}
\end{figure}
Megfordítva, ha vasból készült ilyen üreget izzítunk, a kísérletek szerint a lyuk a környező izzó vasnál jóval erősebben világít. A "fekete testsugárzásának" forrásául ezért állandó hőmérsékleten tartható, rendszerint
elektromosan izzított üreg nyílását használják, és emiatt a fekete test sugárzását másképpen "üregsugárzásnak" is hívják. Más példákat is találhatunk környezetünkben, amelyek bár kívülről nem tlínnek abszolút fekete testnek, abszorpcióképességük mégis 1 (minden rájuk eső sugárzást elnyelnek),
és sugárzásuk ugya nolyan, mint a fekete test s ugárzása. Ilyenek a csillagok, például a Nap, és az izzított fémszál, például a villanykörte izzószálja.
Másik érdekes példa a mikrohullámú háttérsugárzás . Ez egy olyan sugárzás, melynek eredete je len könyv határain kívül esik ugyan (az Univerzum fejlődésének korai szakasza), de tulajdonságai a fekete test sugárzásának
tulajdonságaival megegyeznek. 1964-ben Tokióból volt először távolsági élő televíziós olimpiai közvetítés az USA-ba. A nagy távolságról érkező elektromágneses jelek zajszintjének csökkentésén fáradozva A. A. Penzias és R. W.
Wilson fed ezték fel és később azonosították ezt a mikrohullámú háttérsugárzást, ami az űrből minden irányból érkező elektromágneses sugárzás.
Frekvenciaeloszlása alapján egy 2,7 K hőmérsékletű feketetest sugárzásásával egyenértékű. Ez a mikrohullámú háttérsugárzás a kozmológiai elméletek
számára fontos ta pas ztalat , és az univerzum kialakulásá nak egy ik hírnöke. Felfedezésükért Penzias és W ilson később 1978- ban Nobel-díjat kaptak.

\subsubsection{Kirchhoff törvények}
Hogy a sugárzásra vonatkozólag kvantitatív összefüggéseket írhassunk fel, definiáljuk valamely test abszorpcióképességét, $a$-t. Ez a testre eső sugárzás energiájának az a törtrésze, amelyet a test elnyel, te hát nem ereszt át és nem
ver vissza. Hasonló módon definiálható a $d$ áteresztőképesség és az $r$ reflektálóképesség. Ezekre fennáll, hogy a $a + d + r = 1$, a bejövő energia csak ezen
három módon tud átalakulni, és az összeg megmarad. Továbbá bevezethetjük a test emisszióképességét, ez a test felületének 1 m$^2$-es darabja által
1 s alatt a felületre merőleges egységnyi térszögben kisugárzott energiát jelenti.

Mind $a$ mind pedig a $e$ a $T$ hőmérsékleten és a $\lambda$ hullámhosszon kívül nagymértékben függ a test különböző sajátságai tól, pl. a felület érdességétől, sötétségétől. Pl. már igen vékony karomréteg a látható fényt csaknem teljesen elnyeli, vagyis közelítőleg $a = 1$. Azt a testet, mely a ráeső bármely
hullámhosszúságú sugárzást teljesen elnyeli, abszolút fekete testnek nevezzük. Ennek abszorpcióképessége tehát: $a = 1$, emisszióképességét pedig jelöljük $E$-vel.

Bármennyire is mások $e$ és $a$ értékei az egyes testeknél, a tapasztalat szerint az $e/a$ viszony minden testnél ugyanaz, csak $\lambda$-nak és $T$-nek a függvénye. Ez az $E(\lambda, T)$ függvény az abszolút fekete test emisszióképessége:
\begin{equation}
    \frac{e}{a} = \frac{e_1}{a_1} = \frac{e_2}{a_2} = \ldots = \frac{e_{\text{feketetest}}}{1} = E(\lambda, T)
\end{equation}
Ez a fontos összefüggés Kirchhofi törvénye (1860), amely termodinamikai úton levezethető. Kiemeljük, hogy a törvény- mint az az $e/a= E(\lambda, T)$
írásmódból is kitűnik - minden egyes hullámhosszra külön-külön is vonatkozik. Kirchhoff törvényének
\begin{equation}
    e(\lambda, T) = a(\lambda) \cdot E(\lambda, T)
\end{equation}
alakjából, figyelembe véve, hogy $a < 1$, az következik, hogy az abszolút fekete test emisszióképessége bármely más test emisszióképességénél nagyobb,
ugyanazon hőmérsékleten. Ugyancsak a fenti egyenletből következik, hogy bármely test emisszióképességét megkaphatjuk, ha a fekete test emisszióképességét az illető test abszorpcióképességével megszorozzuk.

\subsubsection{A feketetest sugárzás}
Kérdés, hogy a feketetest sugárzása hogyan is függ pontosan a hullámhossztól és a hőmérséklettől. Ezt a függvényt $E(\lambda, T)$-vel jelöljük, amely megadja, hogy az abszolút fekete test 1 m$^2$ felülete
egységnyi idő alatt a $\lambda$ és $\lambda + d\lambda$ hullámhossztartományban mennyi energiát sugároz ki. Ezt a függvényt nevezzük a fekete test sugárzásának spektrumának. A fekete test sugárzásának spektrumát 
kísérletileg sokat vizsgálták, úgy, hogy egy korommal bekent üreges edényt izzítottak fel, és a lyukból kilépő sugárzás hullámhossz-eloszlását mérték egy termoelem segítségével, amelynek feszültsége a sugárzás intenzitásával arányos.
\begin{figure}[H]
    \centering
    \includegraphics[width=0.5\textwidth]{imgs/10-tetel/fekete_test_spektrum.png}
    \caption{A fekete test sugárzásának spektruma különböző hőmérsékleteken}
\end{figure} 
Az ordinátán lévő $E(\lambda, T)$-t úgy kell érteni, hogy $E(\lambda, T) d\lambda$ az 1 m$^2$ felületről 1 s alatt a $\lambda$ és $\lambda + d\lambda$ hullámhossztartományban kisugárzott energia. A görbékből látható, hogy adott hőmérsékleten van egy
maximális intenzitás, amelyhez egy meghatározott hullámhossz tartozik. Ez a hullámhossz a hőmérséklet növekedésével egyre kisebb lesz, azaz a sugárzás maximuma a rövidebb hullámhosszak felé tolódik el. Ezt a jelenséget Wien eltolódási törvényének nevezzük, amely szerint a sugárzás maximumának hullámhossza és a hőmérséklet szorzata állandó:
\begin{equation}
    \lambda_{\text{max}} T = b
\end{equation}
Itt $b = 2,898 \times 10^{-3} \, \mathrm{m \cdot K}$ a Wien-féle állandó. A kísérleti görbék alapján felírható egy integrált összefüggés is, amely megadja az összes hullámhosszon kisugárzott energiát:
\begin{equation}
    I(T) = \int_0^{\infty} E(\lambda, T) d\lambda = \sigma T^4
\end{equation}
Ez a Stefan-Boltzmann törvény, ahol $\sigma = 5,67 \times 10^{-8} \, \mathrm{W/m^2 K^4}$ a Stefan-Boltzmann állandó.

\subsubsection{Rayleigh-Jeans féle sugárzási törvény}
A fekete test sugárzásának spektrumára a klasszikus fizika többféle elméletet is felállított, de ezek közül csak egyetlen volt képes a kísérleti eredmények egy részét megmagyarázni. Ez a Rayleigh-Jeans féle sugárzási törvény, amelyet Lord Rayleigh (1900) és J. H. Jeans (1905) dolgozott ki. Az elmélet a következő feltevéseken alapult:
\begin{itemize}
    \item Az elektromágneses sugárzás a fekete test üregében állóhullámok formájában létezik, amelyek frekvenciája és hullámhossza a szokásos módon kapcsolódik egymáshoz: $\nu = \frac{c}{\lambda}$.
    \item Az állóhullámok energiája a klasszikus kinetikus gázelmélet szerint oszlik el, azaz minden egyes hullámrezgés átlagos energiája $kT$, ahol $k$ a Boltzmann állandó.
\end{itemize}

Egy állóhullámot a következő módon tudunk felírni az üregben:
\begin{equation}
    \psi(x) \sim \sin kx \quad \text{ahol} \quad \psi(0) = \psi(a) = 0 \quad \text{határfeltétel}
\end{equation}
Ebből a határfeltételből adódik, hogy az $k$ hullámszámra a következő feltétel adódik:
\begin{equation}
    ka = n\pi \quad \rightarrow \quad k = \frac{n\pi}{a} = \frac{2 \pi}{\lambda} = 2\pi \frac{\nu}{c} \qquad n = 0, 1, 2, \ldots
\end{equation}
Mivel a feketetest tulajdonképpen egy üregrezonázor, ezért három irányban is felírhatjuk a hullámszám feltételét, így a következő összefüggést kapjuk:
\begin{equation}
    k_xa = n_1 \pi, \quad k_y a = n_2 \pi, \quad k_z a = n_3 \pi
\end{equation}
Innen a teljes hullámszám:
\begin{equation}
    k^2 = \frac{\pi^2}{a^2} (n_1^2 + n_2^2 + n_3^2)^2 = \frac{4 \pi^2 \nu^2}{c^2}
\end{equation}

Ahhoz hogy meg tudjuk határozni a lehetséges módusok $Z$ számát $\nu' < \nu$ esetére, vennünk kell egy gömböt, amelynek sugara $R$ a $n_1$, $n_2$, $n_3$ koordinátarendszerben:
\begin{equation}
    R^2 \equiv n_1^2 + n_2^2 + n_3^2 = \left( \frac{2 \pi \nu}{c} \frac{a}{\pi} \right)^2
\end{equation}
A teljes térfogatban a $\nu$ frekvenciánál kisebb módusok száma $Z$ az $R$ sugarú gömb térfogatának nyolcada, szorozva a kétféle polarizációval:
\begin{equation}
    Z(\nu) = 2 \cdot \frac{1}{8} \cdot \frac{4}{3} \pi R^3 = \frac{4 \pi}{3} \left( \frac{a \nu}{c} \right)^3
\end{equation}
A $\nu$ és $\nu + d\nu$ közötti módusok száma a $ V = a^3$ térfogatban:
\begin{equation}
    dZ(\nu) = 8 \pi \left( \frac{a}{c} \right)^3 \nu^2 d\nu = \frac{8\pi}{c^3} V \nu^2 d\nu
\end{equation}
De mennyi energiát képviselnek ezek az állóhullámok? Tegyük fel, hogy egy módusra jutó átlagos energia $\bar{\epsilon}$. Ekkor a $V$ térfogatban a $\nu$ és $\nu + d\nu$ közötti hullámhosszakhoz tartozó energia egységnyi térfogatra:
\begin{equation}
    dW = \bar{\epsilon} \cdot \frac{dZ(\nu)}{V} = \bar{\epsilon} \cdot \frac{8\pi}{c^3} \nu^2 d\nu \equiv U(\nu, T) d\nu \equiv 
\end{equation}
Tehát az energia:
\begin{equation}
    U(\nu ,T) = \bar{\epsilon} \cdot \frac{8\pi}{c^3} \nu^2
\end{equation}

Az átlagos energiát ami egy módusra (oszcillátor) jut a klasszikus kinetikus gázelmélet alapján meg lehet határozni. Ha van $N_{\text{tot}}$ darab megkülönböztethető objektumunk, és a rendszer teljes energiája $E_{\text{tot}}$, akkor az egy objektumra jutó átlagos energia:
\begin{equation}
    \bar{\epsilon} = \frac{E_{\text{tot}}}{N_{\text{tot}}}
\end{equation}
Tegyük fel továbbá, hogy az $N_n$ objektum az $n$-edik energiaállapotban van, amelynek energiája $\epsilon_n$:
\begin{equation}
    \epsilon_n = n \epsilon_0 \qquad n = 1, 2, \ldots
\end{equation}
Ekkor az $\epsilon_0 \to 0$ határesetben a módus-csoportok között az energia folytonosan változik.

A Bolztmann eloszlás szerint az $n$-edik állapotban lévő objektumok száma:
\begin{equation}
    \frac{N_n}{N_{\text{tot}}} = \frac{e^{-\epsilon_n / k_BT}}{ \sum_{n=0}^{\infty} e^{-\epsilon_n / k_BT}}
\end{equation}
Tehát az átlagos energia ($\beta = 1/k_BT$):
\begin{equation}
    \begin{aligned}
    \bar{\epsilon} &= \frac{E_{\text{tot}}}{N_{\text{tot}}} = \sum_{n} \frac{N_n}{N_{\text{tot}}} \epsilon_n = \frac{\sum_{n=0}^{\infty} \epsilon_n e^{-\beta \epsilon_n}}{\sum_{n=0}^{\infty} e^{-\beta \epsilon_n}} =  \frac{\sum_{n=0}^{\infty} n \epsilon_0 e^{-\beta n \epsilon_0}}{\sum_{n=0}^{\infty} e^{-\beta n \epsilon_0}} = \\
    &= \frac{-\frac{\partial}{\partial \beta }\sum_{n=0}^{\infty} e^{-\beta n \epsilon_0}}{1 + e^{-\beta \epsilon_0} + (e^{-\beta \epsilon_0})^2 + \ldots} = \frac{-\frac{\partial}{\partial \beta }\frac{1}{1 - e^{-\beta \epsilon_0}}}{\frac{1}{1 - e^{-\beta \epsilon_0}}} = \\
    &= \frac{\frac{\epsilon_0 e^{-\beta \epsilon_0}}{(1 - e^{-\beta \epsilon_0})^2}}{\frac{1}{1 - e^{-\beta \epsilon_0}}} = \frac{\epsilon_0 e^{-\epsilon_0 / k_BT}}{1 - e^{-\epsilon_0 / k_BT}} = \\
    &= \frac{\epsilon_0}{\frac{\epsilon_0}{k_BT} + \frac{1}{2} \left( \frac{\epsilon_0}{k_BT} \right)^2 + \ldots} = \frac{k_BT}{1 + \frac{1}{2} \frac{\epsilon_0}{k_BT} + \ldots} \\
    \bar{\epsilon} &= \frac{\epsilon_0}{e^{\epsilon_0 / k_BT} - 1}
    \end{aligned}
\end{equation}
Ez pedig a klasszikus határesetben, amikor $\epsilon_0 \to 0$ (ekvipartíció tétele):
\begin{equation}
    \bar{\epsilon} = k_BT
\end{equation}
Így a Rayleigh-Jeans féle sugárzási törvény:
\begin{equation}
    U(\nu, T) d\nu = \frac{8 \pi}{c^3} k_BT \nu^2 d\nu
\end{equation}
Ez az összefüggés jól leírja a kísérleti görbéket kis frekvenciákon (nagy hullámhosszakon), de nagy frekvenciákon (kis hullámhosszakon) de akad vele egy kis probléma:
\begin{equation}
    I(T) = \int_0^{\infty} U(\nu, T) d\nu = \int_0^{\infty} \frac{8 \pi}{c^3} k_BT \nu^2 d\nu = \infty
\end{equation}
Ez azt jelenti, hogy a Rayleigh-Jeans féle sugárzási törvény szerint a fekete test végtelen mennyiségű energiát sugároz ki, ami nyilvánvalóan ellentmond a kísérleti tapasztalatokkal. Ezt a problémát nevezzük a \textbf{ultraibolya katasztrófának}.

\subsubsection{Wien-féle sugárzási törvény}
Wien hasonlóan a Bolztmann eloszlásból indult ki, de ő azt feltételezte, hogy a hullámhossz függ a molekulák sebességétől. Ebből a feltevésből kiindulva a következő sugárzási törvényt kapta:
\begin{equation}
    U(\nu, T) \sim \nu^3 e^{\frac{- h \nu}{k_BT}}
\end{equation}
Ahol $h$ egy új univerzális állandó, amelyet később Planck állandónak neveztek el (ezt az állandót természetesen Wien eredeti eredménye még nem tartalmazta). Ez a sugárzási törvény jól leírta a kísérleti görbéket nagy frekvenciákon (kis hullámhosszakon), de kis frekvenciákon (nagy hullámhosszakon) nem működött jól.

\subsubsection{Planck-Féle sugárzási törvény}
A XIX. század végén a hőmérsékleti egyensúlyban lévő fekete test sugárzásának leírására két formulát használtak, de egyik sem tudta értelmezni a teljes
hullámhossztartományt. Az egyik a Rayleigh- Jeans-törvény, amely a spektrum hosszúhullámú részében egyezik meg a tapasztalattal, a másik a Wien
által megállapított formula, amely a rövidhull ámú határesetben érvényes.A sugárzást az egész spektrumban helyesen leíró formulát Max Plancknak sikerült felállítani 1900-ban.

Ő azt gondolta, hogy a Rayleigh-Jeans-törvény azért "bukik el" magas frekvenciákon, mert az $\bar{\epsilon}(T)$ esetleg túl erősen van képviselve nagyfrekvenciás tartományokon.
Ezért a klasszikus $\bar{\epsilon} \sim k_BT$ helyett egy $\bar{\epsilon} < k_BT$ egyenlőtlenséget feltételezett nagy $\nu$ tartományban. Ezt könnyű elérni ha feltesszük, hogy az alapállapot:
\begin{equation}
    \epsilon_0 = h \nu
\end{equation}
ahol $h$ egy új univerzális állandó, amelyet később Planck állandónak neveztek el. Ekkor az átlagos energia:
\begin{equation}
    \bar{\epsilon} = \frac{h \nu}{e^{h \nu / k_BT} - 1}
\end{equation}
Így a Planck-féle sugárzási törvény:
\begin{equation}
    \boxed{U(\nu, T) = \frac{8 \pi h \nu^3}{c^3} \frac{1}{e^{h \nu / k_BT} - 1}}
\end{equation}
Ez az eredmény pedig tökéletesen visszaadja a Rayleigh-Jeans törvényt kis $\nu$-re és a Wien törvényt nagy $\nu$-re:
\begin{equation}
    U(\nu, T) \approx \begin{cases}
    \frac{8 \pi}{c^3} k_BT \nu^2 & \text{ha } h \nu \ll k_BT \quad \text{(Rayleigh-Jeans)} \\
    \frac{8 \pi h}{c^3} \nu^3 e^{-h \nu / k_BT} & \text{ha } h \nu \gg k_BT \quad \text{(Wien)}
    \end{cases}
\end{equation}
Mi több, ha kiszámoljuk a teljes energiát, pontosan a Stefan-Boltzmann törvényt kapjuk:
\begin{equation}
    I(T) = \int_0^{\infty} U(\nu, T) d\nu = \frac{8 \pi k_B^4}{15 c^3 h^3} T^4 \int_{0}^{\infty} \frac{x^3}{e^x - 1} dx = \sigma T^4
\end{equation}
ahol $\sigma = \frac{2 \pi^5 k_B^4}{15 c^2 h^3}$ a Stefan-Boltzmann állandó.

Viszont ez az eredmény sokkal messzemenőbb implikációkkal rendelkezik, mint csak a fekete test sugárzásának leírása. Ugyanis, ha feltételezzük, hogy az elektromágneses sugárzásnak van egy diszkrét minimum értéke, $h \nu$, 
És ennek csak az egész számú többszöröseinek értékét heveti fel ($\epsilon_n = n h \nu$), akkor ez azt jelenti, hogy az elektromágneses sugárzás kvantált. 

Plancknak az energia kvantáltságára vonatkozó feltevését később Einstein módosította, magára a sugárzási térre is alkalmazta a fényelektromos-hatás
kisérleti tapasztalatai alapján. Megmutatta, hogy Planck hipotézise helyett helyesebb, ha feltesszük, hogy a sugárzás emissziója és abszorpciója $h\nu$ energiájú adagokban - kvantumokban - történik,
azaz Einstein hipotézise szerint a fény nem folytonos, hanem korpuszkuláris szerkezetű : a sugárzás $E = h\nu$ energiájú kvantumokból áll. (Érdemes
megjegyezni, hogy Newton a fényt korpuszkuláris természetűnek tartotta, szemben Huygens-szel, aki a hullámtermészet alapján magyarázta az optikai jelenségeket. Természetesen az einsteini fénykvantum nem ugyanaz, mint
Newton "fényrészecskéje".) A sugárzási törvény jelentősége messze túlmegy a hőmérsékleti sugárzás területén, hiszen az anyagra és az elektromos töltésre vonatkozó atomisztikus felfogást itt alkalmazták először az energiára is: az oszcillátor energiája
csak diszkrét értékeket vehet fel, és így az csak ugrásszerűen változhat. Ennek a gondolatnak a kifejlődése egész fizikai világszemléletünket lényegesen átalakította, és a kvantumelmélethez vezetett.

\subsection{Rutherford-kísérlet}
\subsubsection{Atommodellek Rutherford Előtt}
A kinetikus gázelmélet alapján ismert térfogatú anyagban leszámolva az atomok számát, meg lehetett határozni az egyes atomok méretét: a gömb alakúnak tekintett atom sugara $\sim 10^{-10}\,\mathrm{m}=1\,\text{\AA}$ érték adódott.
Az atom belső szerkezetére vonatkozólag azonban nem voltak ismereteink, de a kísérletek alapján fel kellett tételezni, hogy az atomok belső szerkezete bonyolult. Felfedezték az izzított fémek termikus elektronemisszióját, a külső feszültséggel előidézett katódsugárzást, gázkisülést, megismerték a villámot, melyek során elektronok léphetnek ki az atomokból.
A gázokban lejátszódó folyamatok is arra mutattak, hogy az atomok belsejében pozitív és negatív elektromos töltés található. Az első felfedezett radioaktív atomok $\alpha$-részecskéket bocsátanak ki. Ez adta kezünkbe az eszközt, melynek felhasználásával széleskörű munka indult az atomok szerkezetének vizsgálatára.

A múlt század végén ismert volt az elektromágneses tér máig is stabilan álló négy pillére, a Maxwell-egyenletek rendszere. Ehhez képest keveset tudtunk az elektromosságot létrehozó töltések tulajdonságairól, így az első elemi részecskék viselkedéséről. Az elektront 1897-ben fedezték fel olyan értelemben, hogy tulajdonságait az akkori műszereknek megfelelő lehető legnagyobb pontossággal meghatározták.
Az $\alpha$-részecskék kibocsátásával járó radioaktivitást 1898-ban fedezték fel, és csak később sikerült az $\alpha$-bomló preparátumból az $\alpha$-részecskéket kisülési csőbe gyűjteni, majd a színképe alapján azonosítani (1909). A hidrogénatom pozitív töltésű ionját már ismer(t)ék a századfordulón, arra azonban, hogy ez a többi atomban alkotóelemként is szerepel, még sokáig nem derült fény.

Thomson az általa felfedezett elektronra alapozta első atommodelljét, a „mazsolás kalácsot” (1900). Ez egy — a kinetikus elmélet alapján — $10^{-10}\,\mathrm{m}$ átmérőjű gömb, mely kívülről semleges. Benne a pozitív töltések egyenletesen vannak elosztva a teljes térfogatban, és ebbe beágyazva pontszerű elektronok helyezkednek el.
Az elektronok vagy nyugalomban vannak, vagy meghatározott pályákon körbe keringenek. Thomson sok számítást végzett az elektronrendszerek stabilitására, és sikerült kimutatnia, hogy a megfelelően elhelyezett elektronrendszer tagjainak sugárzása leronthatja egymást, így gyakorlatilag stabil atomot eredményez. A század elején, 1903-ban, P. Lenard a katódsugarakkal végzett méréseket.
Az ilyenkor előállított gyors elektronok vékony fémfóliákon és gázokon való elnyelődését és szóródását vizsgálta. Azt találta, hogy az elektronok nagyon sok atomon áthaladva sem nyelődnek el, áthatolnak lényeges irányváltoztatás nélkül. Ebből arra következtetett, hogy az „áthatolhatatlan atom” nem tölti ki a rendelkezésére álló teret, az atom nagy része üres. Lenard úgy gondolta, hogy az atom kis méretű pozitív és negatív töltésekből, valamint azok erőteréből áll.
Ez a modell nem illeszkedett az akkori elméleti számításokhoz, nevezetes ellentétben a klasszikus elektrodinamika jóslataival, így Lenard modellje nem lett olyan elfogadott, mint Thomsoné. De hogyan helyezkednek el a pozitív és negatív töltések az atomban valójában? Angliában a század első éveiben Rutherford és munkatársai, Geiger és Marsden, az akkor nem régen felfedezett $\alpha$-sugárzás részecskéivel bombáztak vékony aranyfóliákat. Meg akarták mutatni, hogy a „mazsolás kalácsban” hogyan térülnek el a kétféle pozitív $\alpha$-magok az elektrodinamika törvényeinek megfelelően.

\subsubsection{Kisérleti elrendezés}
Az Ernest Rutherford vezette kísérletekben (1906--1911) párhuzamos
$\alpha$-részecskenyalábot bocsátottak egy céltárgyra, és regisztrálták
a részecskék haladási irány szerinti eltérülésének eloszlását.
A céltárgy általában a vizsgált atomokból, az első kísérletekben aranyból
álló vékony fólia vagy gázréteg volt, melyben a többszörös szórás
valószínűsége igen csekély, így a kísérletek értelmezése egyszerűsödött.
A párhuzamos $\alpha$-részecske-nyaláb anyagrétegen áthaladva
szétszóródik.
Az $\alpha$-részecskék többsége kis szögű, $2$--$3^\circ$-os szórást szenved.
A szórási szögek eloszlása közelítőleg Gauss-eloszlást követ a
„mazsolás kalács”-modellben.
Ezeken a kis szöggel történő szóródásokon kívül azonban, mint Rutherford
munkatársai, Geiger és Marsden megmutatták, az $\alpha$-részecskék kb.
$\tfrac{1}{8000}$-ad része $90^\circ$-ot is meghaladó nagy szögben
szóródik vissza.
A nagyszögű szóródások száma jelentősen meghaladja a Gauss-eloszlás
által megengedett, azaz a véletlenből fakadó mértéket
(statisztikusan sok kis szög felhalmozódása).

A szóródásnak az eltérülés irányától való függését vizsgáló készülékük
vázlatát az 5.1. ábra mutatja.
A berendezés egy vastag falú hengeres edény, amely csiszoltaban
forgatható módon volt elhelyezve.
Az M leolvasó mikroszkóp és az ehhez rögzített S cink-szulfid ernyő a
dobozszal együtt forgott.
Az $\alpha$-forrás és az F szórólemez mereven össze van építve,
maga a szórólemez a doboz forgási tengelyében van.
A cink-szulfid ernyőn minden egyes nagy energiájú $\alpha$-részecske
becsapódás sötétben jól felismerhető fényfelvillanást (ún. szcintillációt)
okoz.
\begin{figure}[H]
    \centering
    \includegraphics[width=0.6\textwidth]{imgs/10-tetel/rutherford.png}
    \caption{Rutherford-kísérlet vázlata ($R$: $\alpha$-forrás, $F$: szórófólia, $S$: cink-szulfid ernyő, $M$: mikroszkóp)}
\end{figure}
Valamely irányban szóródó $\alpha$-részecskék számát úgy határozták meg,
hogy a megfelelő szög beállítása után leszámlálták a felvillanásokat a
szcintilláló ernyőn.
Az $\alpha$-részecskék forrása vékony falú, radonnal töltött cső volt,
a sugárnyalábot egy keskeny diafragma szűkítette le.
A mérések $30^\circ$--$150^\circ$ közötti szögintervallumra terjedtek ki.
A szög függvényében ábrázolták az észlelések számát, vizsgálva, hogy
milyen függvény szerint változtak.
Vizsgálták a szórás függését a $\theta$ szögtől kívül a bombázó részecskék
sebességétől, valamint a céltárgyban levő szórócentrumok számától is.
A szögeloszlásadatok nem követték a „mazsolás kalács” által jósolt
Gauss-eloszlást, hanem
$\displaystyle \frac{1}{\sin^4\frac{\vartheta}{2}}$ függvény szerint a nagy szögekben
egyre kevesebb $\alpha$-részecske szóródott, de ez az eloszlás megengedti
a nagy szögű visszaszórást is.
Az $\alpha$-részecskék sebességét úgy lehetett változtatni, hogy más
energiájú $\alpha$-sugarzó izotópot helyeztek a radon helyére (pl. polónium,
rádium).
A kísérleti eredmény az lett, hogy minél nagyobb az $\alpha$-részecske
energiája, annál kisebb a szóródás valószínűsége egy adott szögben.

Manapság iskolai körülmények között elérhető, hogy a Rutherford-kísérletet
elvégezzük (és az atommag létezését megmutassuk), természetesen nem
cink-szulfid ernyővel.
Az $\alpha$-részecskék forrását, az ólomkollimátort és az aranyfóliát
ugyanúgy alkalmazva, de a cink-szulfid ernyő helyett egy félkör alakban meghajlított, ún. nyomdetektort (ami egy speciális fotólemez)
elhelyezve az a-részecskék irány szerinti eltérülése vizsgálható. A kísérlet
legnagyobb nehézsége, hogy vákuumba kell helyezni a berendezést, mert
a levegő atommagjai szintén eltérítik az a-részecskéket, és ez elmossa az
aranyfólia hatását. A nyomdetektorlemezt a több órás besugárzás után erős
$NaOH$-oldattal maratni kell, így minden egyes a-részecske becsapódása helyén egy kis lyuk keletkezik, melyek száma mikroszkóp alatt megszámolható.
Összegezve a Rutherford-féle szóráskísérletek legfontosabb kisérleti tapasztalatait, a következőket mondhatjuk:
\begin{enumerate}
    \item Az $\alpha$-részecskék $\vartheta > 90^\circ$-os, nagy szögű visszaszórását lehet megfigyelni.
    \item A beütések szögeloszlása nem Gauss-eloszlást, hanem $\frac{1}{\sin^4\frac{\vartheta}{2}}$ függvény szerinti eloszlást mutat, egy adott detektorral végignézve a $30\degree- 150\degree$ eltérülési szögtartományt
\end{enumerate}

\subsubsection{Rutherford atommodellje}
Rutherford arra a következtetésre jutott, hogy ilyen nagy szögű eltérések
csak abban az esetben lehetségesek, ha az atomban a tömeg igen kis
térfogatra koncentrálódik, melynek pozitív töltése van, és amelynek
sugara jóval kisebb az atom méreténél, ehhez képest az elektrodinamikai
modellben pontszerűnek tekinthető.
Ez Rutherford nagy felfedezésének lényege, hogy az atomon belül a tömeg
nagy része (kb. $99{,}9\%$-a) egy pici pozitív töltésű térrészben
sűrűsödik (ezek a visszaszóró centrumok az atomok „magjai”,
melyek töltését magtöltésnek nevezik), a maradék térrész pedig csak
lazán van kitöltve ezek szerint negatív töltésű részecskékkel
(az atom egésze semleges), kézenfekvően elektronokkal.
Ez a kísérlet tekinthető az atommag felfedezésének.
A Rutherford által javasolt atommodellben az atom felépítése a
bolygórendszerhez hasonlít.
Középen van a kisméretű pozitív töltésű atommag, amelyben az atom
tömegének túlnyomó része koncentrálódik, míg a mag körül zárt pályán
negatív elektronok „keringenek” az elektromágneses kölcsönhatás ismert
törvényszerűsége, a Coulomb-erő alapján.
A Coulomb-erő alakilag teljesen megegyezik a tömegvonzás newtoni
erőtörvényével, innen jön a bolygórendszer-analógia.

Ha feltételezzük, hogy a besugárzott $\alpha$-részecskék nem érnek
hozzá az atommaghoz (a következő pontban be is látjuk), ebből már egy
becslést lehet adni az atommag méretére, centrális ütközést feltételezve.
Az $\alpha$-részecske körülbelül $8\ \mathrm{MeV}$ mozgási energiája a
céltárgyhoz legközelebbi pontban teljesen átalakul elektromos helyzeti
energiává.
Az
\[
E = k\,\frac{z_1 z_2 e^2}{r}
\]
egyenlet alapján a mag sugarának egy felső korlátja
$\sim 100\cdot 10^{-15}\ \mathrm{m} = 100\ \mathrm{fm}$-nek adódik.
Az atommagok méretének pontosabb mérései alapján a sugaruk
$1$--$10\ \mathrm{fm}$ nagyságrendbe esik, tehát tíz- vagy
százezerszer kisebbek az atom méreténél.

Egy modell igazolásához természetesen elméleti számolásokat kell végezni,
amelyek a kísérlettel megegyező eredményre vezetnek.
Rutherford a kísérletek alapján megalkotott atommodell segítségével
kidolgozta az $\alpha$-részecskék szórásának elméletét, mellyel a
kísérleti adatokat sikerrel írta le, és amelynek végeredménye a
Rutherford-féle szórási formula nevet kapta.

\subsubsection{Hatáskeresztmetszet}
Az atomfizikában a reakciók valószínűségeit a hatáskeresztmetszet
fogalmának segítségével jellemezhetjük.
Itt a Rutherford-formula éppen az $\alpha$-részecskék aranyatommagokon
történő tisztán elektromágneses szóródásának hatáskeresztmetszetét adja
meg, ezért először néhány szót ejtünk a hatáskeresztmetszet fogalmáról
általában.

\subsubsection*{A hatáskeresztmetszet fogalma}
A hatáskeresztmetszet egy adott atom- vagy magfizikai reakció
lejátszódásának valószínűségére jellemző fogalom.
Megmutatja például, hogy a lassú neutronok mennyire szeretnek befogódni
egyes atommagokba, vagy mekkora a valószínűsége annak, hogy egy
atommagot protonokkal bombázva egy adott térszögben deutérium magok
lépnek ki.
Hasonló módon megmutatja, hogy egy $\alpha$-részecskenyaláb részecskéi
milyen valószínűséggel fognak (pl.) $10^\circ$-os szögben eltérülni.
Ezekben a példákban az a közös, hogy egy adott részecske párhuzamos
nyalábja okozza a reakciót.
Legyen a nyaláb keresztmetszete $F$ és tekintsük homogénnek, valamint
először legyen csak egy darab bombázott atom vagy atommag jelen.
Ekkor, ha adott $\Delta t$ idő alatt $N$ részecske érkezik a nyalábbal
és ebből $N_r^{(1)}$ vesz részt a reakcióban, akkor a reakció
valószínűsége egy elemi reakcióban
\[
p=\frac{N_r^{(1)}}{N}.
\]
Ez megadható a reakcióban részt vevő részecskéknek a nyalábban
elfoglalt felületének ($\sigma$) és a teljes nyalábkeresztmetszetnek
($F$) hányadosaként is:
\[
p=\frac{N_r^{(1)}}{N}=\frac{\sigma}{F}.
\]

Természetesen sok atommagot bombázva a reakciók átlagos száma arányos
a bombázott magok számával is ($N_c$) (c = céltárgy).
Az előző egyenletet átrendezve és általánosítva kapjuk:
\[
N_r = p\cdot N\,N_c = \frac{\sigma}{F }N\,N_c
= \sigma\frac{N}{F \Delta t}\,N_c \Delta t = \sigma j N_c \Delta t.
\]

Itt $j$ a részecske-áramsűrűség, a nyalábban a másodpercenként és
felületegységenként áthaladó részecskék számát adja meg.
A bombázó részecskék közül annyi darab vesz részt a reakcióban (egy adott atommaggal történő
reakcióban), ahány a „hatásos” $\sigma$ nagyságú keresztmetszetbe befér.
Ez a $\sigma$ a hatáskeresztmetszet, és az alábbi egyenlet a definíciója:
\[
\frac{dN_r}{dt}=\dot{N}_r=\sigma j N_c.
\]

Látható, hogy a valószínűség geometriai értelmezésével jutunk a
hatáskeresztmetszet szemléletes értelméhez, a nyalábnak ezen felületű
része fog részt venni a reakcióban, vagy adott irányba szóródni.
(Nem kell mindig összefüggő tartományt képzelni a hatásos keresztmetszethez.)
Az érdekes csak az, hogy nem a $p$ valószínűség, hanem a „hatásos
keresztmetszet” a céltárgy magra általánosan jellemző állandó.

Ha egységnyi térfogatban a bombázott részecskék $N_c/V=n_c$ darab
hatásos keresztmetszeteit összeadjuk (az átfedés valószínűsége
elhanyagolható), akkor a céltárgy teljes, ún. makroszkopikus
hatáskeresztmetszetét kapjuk.
Ennek jele a görög nagy szigma ($\Sigma$).
$\Sigma=\sigma\cdot n_c$, ezzel:
\[
\dot{n}_r=\Sigma j.
\]

A makroszkopikus hatáskeresztmetszet egy adott kísérleti elrendezésre
jellemző fogalom.
A bombázott magok száma ugyan fogy az időben (ha a reakciókban
átalakulnak), és így a $\Sigma$ is csökken, de ez az esetek döntő részében
nem észlelhető egy kísérletben.

További formulák a magreakciók számának meghatározására

a) A $j$ itt részecske-áramsűrűség, nem az elektromos áramsűrűség.
Átalakíthatjuk a formulát a részecske-áramerősségre is, $j=I/F$ alapján,
ilyenkor az $N_c=n_c V=n_c F x$ helyettesítéssel az
$N_r=\sigma j N_c \Delta t$ adódik, amiben a besugárzott anyagdarab
vastagságát $x$ játszik szerepet.
Az $n_c$ a céltárgy részecskesűrűségét jelenti.
Így jutunk az összefoglaló egyenletekhez, $N_r/\Delta t=\dot{N}_r$:

(5.1)
\[
\dot{N}_r=\sigma j N_c=\sigma I n_c x.
\]

Az eddig tárgyalt formulák olyan esetben hasznosak, amikor egy
irányított nyaláb érkezik egy anyagdarabra, általában vékony vagy vastag
fóliára.
Ez az eset a mag- és részecskefizikai kísérletekben általános.

b) Előfordul azonban, hogy a reakciót kiváltó részecskék nem irányítottan
érkeznek, hanem sebességvektoruk a tér minden irányában mutathat.
Ilyenkor a $j$ áram­sűrűség nem egyértelmű, helyette a $\Phi$ fluxus
használandó, ami megadja az egységnyi felületen időegység alatt
bármilyen irányban átmenő részecskék számát.
Tipikus példa erre egy atomreaktor aktív zónája, ahol pl. a neutronok
összevissza mozoghatnak.
Ekkor:
\[
\dot{N}_r=\sigma \Phi N_c.
\]
Ebben az esetben az az előfordulhat, hogy a bombázó részecskék nem
monoenergetikusak, és a hatáskeresztmetszet általában függ a bombázó
részecske energiájától.
Ilyenkor az adott sebességű neutronok (bombázó részecskék) fluxusát az
energia függvényében kell ismerni:
\[
\phi(E)=\frac{d\Phi(E)}{dE},
\]
ez azt jelenti, hogy $\phi(E)$ a $[E,dE]$ energiaintervallumba eső része
a teljes $\Phi$ fluxusnak.
Ilyenkor a keletkező részecskék számát külön-külön a bombázó energiákra
összegezni kell:
\[
\dot{N}_r=\int \sigma(E)\,\phi(E)\,dE\cdot N_c.
\]

c) Geometriáját tekintve még bonyolultabb a modell, ha a céltárgyak
is mozognak.
(Reaktorban az urán és a moderátor atommagok csak szobahőmérsékletű
hőmozgást végeznek, vagyis a bombázó magokhoz képest alig mozognak!)
Ez valósul meg például egy csillag belsejében, amikor a hidrogén- és
héliummagok is több millió fokos gázként viselkednek.
Ilyenkor a részecskesűrűségekkel kell dolgoznunk.
Adott $V$ térfogatban $N_c=n_c\cdot V$ és $N=n\cdot V$.
A bombázónak tekintett magok áramsűrűsége $j=n\cdot v$ alapján számolható
(lásd klasszikus fizika).
Ezekkel
\[
\dot{N}_r=\sigma j N_c=\sigma n v\, n_c V,
\]
ha $V$-vel mindkét oldalt elosztjuk, a térfogategységenkénti reakciók
számát kapjuk meg:
\[
\dot{n}_r=\sigma v\; n_c n.
\]

\subsubsection*{Differenciális hatáskeresztmetszet}

Ha a reakcióban keletkező vagy szóródó részecskék reakció utáni
tulajdonságaiból megkötünk néhányat, például irányát vagy energiáját
külön meghatározzuk, akkor differenciális hatáskeresztmetszetről
beszélünk.
A Rutherford-szórás esetében a szög (térszög) szerinti differenciális
hatáskeresztmetszetről van szó.
Ilyenkor $N_r$ az adott irány körüli egységnyi térszögbe szóródott
részecskék számát adja meg, és $\sigma$ helyett
$\dfrac{d\sigma}{d\Omega}$ jelölést használunk.
$\dfrac{d\sigma}{d\Omega}(\vartheta)$ szemléletesen megadja, hogy
egységnyi idő alatt egy darab céltárgy-atom (mag) a bejövő részecskenyaláb
felületének mekkora részén érkező részecskéket irányítja az adott
irányban lévő egységnyi térszögű detektorba.
A kérdés az, hogy ez a valószínűség hogyan függ az eltérülés szögétől.

\subsubsection{A Rutherford-féle szórási formula}
Rutherford elméletében levezette, hogy az $\alpha$-részecskék egy fix
töltött szórócentrumon történő szóródásának differenciális
hatáskeresztmetszete hogyan függ az eltérülés $\vartheta$ szögétől
\[
\left(\frac{\partial\sigma}{\partial\Omega}(\vartheta)\sim\frac{1}{\sin^4\!\tfrac{\vartheta}{2}}\right).
\]
A kísérleteiben ugyanezt szögfüggést tapasztalta, ez adja modelljének
meggyőző kísérleti alapját.

A modell alapvető feltétele, hogy a pontszerűnek tekintett atommag és az
$\alpha$-részecske közötti kölcsönhatás a Coulomb-törvényt követi, ahogy az
töltött részecskék viselkedése során megszoktuk, de most minden eddiginél
kisebb távolságokra alkalmazzuk.
A centrális $1/r^2$ alakú erők leírásával a mechanikában már találkoztunk
az égitestek mozgása kapcsán.
Az egyenletek teljesen azonos alakúak a két esetben.
A perdület megmaradása miatt síkbeli problémát polárkoordinátákkal
felírva, az energia megmaradását figyelembe véve, a mozgásegyenleteket
megoldva, azt az eredményt kapjuk, hogy a testek pályája egy kúpszelet
\[
r=\frac{p}{1-e\cos\vartheta}.
\]
Ez negatív összenergiánál ellipszis, pozitív összenergiánál (nem kötött
rendszer esetén) hiperbola.
[Budo: Mechanika, 95. oldal.]
A belőtt $\alpha$-részecske aszimptotikusan egy adott irányban eltérül,
ilyenkor a hiperbola esete valósul meg, és az eltérülés szöget a két
aszimptota (bemenő és kifutó) bezárt szöge adja meg.
Ha a céltárgy-mag tömegéhez képest az $\alpha$-részecske tömegét
elhanyagoljuk (vagy a redukált tömeggel
$\mu=\dfrac{m_1 m_2}{m_1+m_2}$ számolunk), akkor a mozgásegyenletek
megoldása után a pályagörbe alakját ismerve a következő formula adódik
az eltérülés szögére:
\[
\tan\frac{\vartheta}{2}
= k\,\frac{Z_1 Z_2 e^2}{m_\alpha v_\alpha^2\, b}.
\]
Itt $m_\alpha$, $Z_\alpha$ és $v_\alpha$ az $\alpha$-részecske tömege,
rendszáma ($Z_\alpha=2$) és kezdeti sebessége, $e$ az elemi töltés,
$b$ az ütközési paraméter;
ebben a távolságban haladna el az $\alpha$-részecske, ha nem lenne
elektromos taszítás a céltárgy-mag és közt(e),
$Z_1$ a bombázott atom rendszáma, $k=\dfrac{1}{4\pi\varepsilon_0}$.
(Ha az $\alpha$-részecske tömegének a helyére a rendszer redukált
tömegét írjuk, akkor figyelembe vettük a szórócentrum véges tömegét is.)

A formulában szereplő $b$ értékét nem tudjuk mérni, ez minden ütközésnél
más és más értéket felvevő valószínűségi változó.
A kísérletekben az eltérülés $\vartheta$ szögét tudjuk csak
meghatározni.
Annak érdekében, hogy ennek eloszlását leíró formulához jussunk,
statisztikus megfontolásokat kell tennünk.

A $\vartheta$ szóródási szög és a $b$ ütközési paraméter egymással az
egyenletben meghatározott kapcsolatban vannak, ezért azok a részecskék
szóródnak $(\vartheta,\vartheta+d\vartheta)$ szögben, melyeknek az
ütközési paramétere $(b-d b,b)$ intervallumba esett,
azaz a beeső nyaláb köré írt $b$ sugarú, $db$ vastagságú körgyűrűn
haladtak át beeséskor.
A bejövő részecskék közül csak azok térülnek el $\vartheta$ szögben,
melyek ezen „hatásos” gyűrűn mentek keresztül.
Avagy: az $F$ teljes keresztmetszet csak azon hányada szóródik
$\vartheta$ szögben, amely ezen a gyűrűn átmegy:
\[
\Delta\sigma=2\pi b\,db.
\]
Ezen hatásos keresztmetszet azonban nem egységnyi térszögbe irányuló
szórásokat vesz csak számba, hanem hengerszimmetrikusan mindet.
El kell osztani még a
\[
\Delta\Omega=2\pi\sin\vartheta\,d\vartheta
\]
térszöggel, ami a térszög $A/r^2$ definíciója alapján adódik:
a gyűrű területe $2\pi r\sin\vartheta\,dr$,
\begin{figure}[H]
    \centering
    \includegraphics[width=0.5\textwidth]{imgs/10-tetel/rutherford2.png}
    \caption{Rutherford-szórás ütközési paramétere és eltérülési szöge}
\end{figure}
A szélessége $dr=r\,d\vartheta$, minden pontja a szórócentrumtól $r$
távolságban van ($r\gg b$).
Így a következő formula adódik:
\[
\frac{d\sigma}{d\Omega}=\frac{b}{\sin\vartheta}\frac{db}{d\vartheta}.
\]

Az $\tan \frac{\vartheta}{2}$ képlet alapján, ha $b$-nek a $\vartheta$-függését kifejezzük,
majd deriváljuk, akkor minden szükséges függvény rendelkezésre áll:
\[
\frac{\partial\sigma}{\partial\Omega}_{\text{Rutherford}}(\vartheta)
=\left(\frac{k e^2 Z_1 Z_\alpha}{2 m_\alpha v_\alpha^2}\right)^2
\frac{1}{\sin^4\!\tfrac{\vartheta}{2}}.
\]

Ez Rutherford formulája az $\alpha$-részecskék szóródására.
A kísérleti eredményekkel való összehasonlítás érdekében a hatáskeresztmetszetből
ki kell számolni, mennyi beütést várunk egy adott kísérlet geometriai
viszonyai között.
A makroszkopikus hatáskeresztmetszet ilyenkor:
\[
\frac{d\Sigma}{d\Omega}=n\frac{d\sigma}{d\Omega},
\]
ahol $n$ a szóró magok koncentrációja.
Jelöljük a céltárgyra másodpercenként és felületegységenként beeső
részecskék teljes számát $j$-vel (részecske-áramsűrűség),
és tételezzük fel, hogy a céltárgy vékony, azaz a kétszeres és többszörös
szórás valószínűsége elhanyagolható.
Ekkor a $\vartheta$ körüli $\Delta\Omega$ térszögbe szóródó részecskék
száma másodpercenként:
\[
\frac{\Delta N_r}{\Delta t}=j\,N\,\frac{d\sigma}{d\Omega_{\text{Rutherford}}}(\vartheta)\,\Delta\Omega.
\]

A Rutherford-formulát azóta az újabb és újabb technikai eszközökkel
gondos kísérleti vizsgálatnak vetették alá, a mérési eredmények minden
esetben jól egyeztek a számítottakkal, amíg a bombázó részecske nem
érintkezik a szórócentrummal.

\subsubsection{A Rutherford-kísérlet következményei}
A Rutherford-formula alapján meg lehet határozni az atommag töltését is,
hiszen a szóródás hatáskeresztmetszete függ $Z^2$-től.
A visszaszórási kísérletben adott aktivitású forrással és geometriai
elrendezés mellett, adott szögben, de a céltárgyat cserélve Chadwick
és munkatársai azt az eredményt kapták (1920), hogy az adott térszögbe
visszaszórt $\alpha$-részecskék számából meghatározott magtöltés
$1$--$2\%$-ra (a kísérlet pontossága is ennyi volt) megegyezik
az elemi töltésnek és az atom rendszámának szorzatával.
(Az eltérített és a beeső $\alpha$-k számát egyidejűen kellett mérni a
magtöltés számértékének meghatározásához.)
A rendszámot akkoriban a Mengyelejev-féle periódusos rendszerben elfoglalt
helyével adták meg.
Miután kiderült, hogy a magtöltés éppen a rendszámnak és az elemi
töltésnek a szorzata, ebből következik, hogy $Z$ darab elektron kering
a mag körül (az atom semlegessége miatt).
A rendszám ezen hármas jelentése (periódusos rendszeri sorszám,
magtöltés és elektronok száma) csak az atommag felfedezése után vált
kézenfekvővé (Van Den Broek, 1913), majd néhány évvel később vált
kísérlettel megerősített ténnyé.

Természetesnek hangzik, de nem egyértelmű, hogy a Coulomb-törvény atomi
méretű töltött részecskék között is érvényes.
Coulomb makroszkopikus ionizált anyagdarabok között tudta mérni a
fellépő elektromos erőt!
A szórási kísérletek erre is felvilágosítást adtak.
Ezek szerint a szórás a Coulomb-féle erőtörvénynek megfelelően megy végbe
mindaddig, amíg a szórt részecskék nem közelítik meg
$10^{-14}\ \mathrm{m}$-nél (azaz $10\ \mathrm{fm}$-nél) jobban az atommagot.
Ez éppen a magok méretének nagyságrendje.
Egyértelműnek tűnik, hogy ennél kisebb távolságon belül azért nem
érvényes a Coulomb-erőtörvény, mert az $\alpha$-részecske hozzáér az
atommaghoz, és más erők is beleszólnak a szórásba.
Még kisebb távolságokon már az elektromos töltés értelmezésével is
gondosan kell bánni, mert kvantumos folyamatok miatt virtuális
elektron-pozitron párok keletkezhetnek a töltések körül, befolyásolva
az elektromos folyamatokat.

A Rutherford-szórás ma már anyagtudományi célú alkalmazott magfizikai
technika, segítségével kristályok felületén lévő szennyező atomok igen
kis koncentrációja is meghatározható.
Azért csak a felületen, mert a néhány megaelektronvolt energiájú
$\alpha$-részecskék nem tudnak néhány mikrométernél jobban behatolni
a szilárd testekbe, ugyanis lelassulnak és megállnak.

Érdemes megjegyezni, hogy Rutherford klasszikus kísérlete óta minden
tisztán elektromágneses atommagon történő szórást Rutherford-szórásnak
hívunk.
Az így felgyorsított elektronok szóródása is a fenti formulát követi,
egészen addig, amíg az elektronok relativisztikussá válnak, amikor is
a Rutherford-formulához korrekciós tagok szükségesek (Mott-szórás).
Másfajta korrekciót jelent az, hogy az atommagok csak egy bizonyos
távolságból tűnnek pontszerűnek.
Az említett gyors elektronok képesek az atommag belsejét is feltérképezni
és az atommagon belüli töltéseloszlásról informálni.
Ha a szórócentrum nem pontszerű, akkor a Rutherford-formulát egy
alakfaktorral megszorozva kapjuk a kísérletekkel egyező és az elméletbe
jól beleillő formulát.

Rutherford atommodelljét gyorsan elfogadta a fizikus közvélemény,
de volt egy gyenge pontja, nem tudta megmagyarázni, hogy miért nem
sugároznak az atommag körül mozgó elektronok.
A periodikus pályákon mindenképpen kell sugárirányú gyorsulásnak lenni,
akkor pedig a klasszikus elektrodinamika szerint energiát kell sugározniuk,
így az elektronok idővel beleesnek a magba.
Ilyet azonban a kísérletekben nem tapasztaltak, ez a század eleji fizika
nyitott kérdése maradt még egy ideig.

\subsection{Millikan-kísérlet}
Az elemi töltést először J. J. Thomson próbálta meg kimérni 1900-ban, amikor az H. A. WIlson által 1895-ben feltalált cseppekondenzációs kamrát használta fel arra, hogy a katódsugárcsőből kilépő elektronok által ionizált levegőben apró vízcseppeket hozzon létre.
Ez a ködkamra jelensége, am i annyit je lent, hogy túltelített gőzben mozgó elektromos töltéseken vízgőz csapódik le, és vízcseppecskék ala.kulnak ki.
Thomson egy gó'zzel telt tartályban röntgensugarakkal keltett ionokat, és a tartály térfogatának hirtelen megnövelésénél bekövetkező lehűléssel tette
a gőzt túltelítetté. A keletkező vícseppecskék egy lemezkére kicsapódtak és azt töltötték fel. Thomsan az elemi töltés nagyságát, $e$-t, a cseppek száma
és az össztöltés hányadosából határozta meg. Ebből az elektron tömegére is megkapta az első eredményeket, és az a hidrogénatom tömegénél három nagyságrenddel kisebbnek adódott.

Thomson módszerét először H. A Wilson javította tovább 1906-ban, majd R. A. Millikan 1911-ben. Millikan amerikai fizikus annyira tökéletesítette a módszert, hogy sikerült ezzel
a mérésével a ma más módszerekkel elérhető pontosságot megközelítenie.
\begin{figure}[H]
    \centering
    \includegraphics[width=0.6\textwidth]{imgs/10-tetel/millikan.png}
    \caption{Millikan cseppkísérlete}
\end{figure}
A kísérlet elve a következő: vízszintes helyzetű síkkondenzátor lernezei közé elektromosan töltött olajcseppecskéket juttatott be. (Más folyadék
cseppjeit vagy szilárd részecskék porát is lehet alkalmazni.) Feszültséget kapcsolva a lemezekre, egy kiszemelt, elektromosan töltött csepp a tér hatására
gyorsulni kezd, s ahogy gyorsul, a levegőmolekulákkal való ütközések miatt bekövetkező belső súrlódás egyre nagyobb fékezőerőt ($F_f$) hoz létre. Végül
Vizsgáljuk meg, milyen erők hathatnak a töltött cseppre, ha be van
kapcsolva az elektrosos tér a kondenzátor lemezei között!
A nehézségi erőn és a felhajtóerőn kívül hat az elektrostatikus erő és a
közegellenállás.
Ha a molekulák szabad úthossza kisebb a cseppek méreténél, akkor a
közeg fékezőereje folyamatosnak (időben állandónak) tekinthető, és
Stokes törvénye értelmében az erő a v sebességgel arányos:

\begin{equation}
F_f = 6\pi\eta r v,
\end{equation}

ahol $\eta$ a belső súrlódási együttható, $r$ az olajcsepp sugara.
Mivel a cseppre állandó gyorsítóerő hat, de a sebességgel arányos
fékezőerő lassítja, ezért egy bizonyos $v_0$ sebességnél az erők
kiegyenlítik egymást, és a sebesség állandó marad.
Legyen a felfelé mutató irány a pozitív, a kondenzátor polaritása olyan,
amely felfelé gyorsítja az elektront.
Ekkor felírva az elektron mozgásegyenletét elektromos térben kapjuk:
\begin{equation}
qE - \frac{4\pi}{3}r^3\sigma g + \frac{4\pi}{3}r^3\varrho g - 6\pi\eta r v_0 = 0
\end{equation}
Itt $q$ a csepp töltése, $E$ az elektromos térerősség, $\varrho$ a levegő
sűrűsége, $\sigma$ az olaj sűrűsége.
Az egyenlet bal oldalának első tagja az elektrostatikus erő, második
tagja a nehézségi erő, harmadik a felhajtóerő és a negyedik a
közegellenállás.
$v_0$-t a mikroszkóp látóterében való megfigyeléssel lehet megmérni.
A csepp sugarát, $r$-t, nem lehet így megmérni, a mikroszkópban
az ilyen kis csepp csak fénylő pontnak látszik.

A sugár meghatározására lehetőséget ad az, ha a sebességet egy másik
E térsség mellett is megmérjük, hiszen az egyenletekben két ismeretlenünk
van, így két egyenlet szükséges.
A számolás szempontjából előnyös, ha az egyik térerősség értékét a 0,
azaz kikapcsolt feszültség mellett is meghatározhatjuk az egyensúlyi
sebességet, ezt nevezzük most $v_1$-nek.
Az előző helyzethez képest a mozgásegyenlet csak annyiban változik,
hogy $E = 0$ és az $F_f$ előjele megfordul, mert így a csepp már lefelé esik.
Ekkor szerencsére a másik ismeretlen ($q$) az $E$-vel együtt eltűnik.
Ilyenkor csak $r$ az ismeretlen, ebből az egy mérésből meghatározható.
Az elektron mozgásegyenlete elektromos tér nélkül:
\begin{equation}
-\frac{4\pi}{3}r^3\sigma g + \frac{4\pi}{3}r^3\varrho g + 6\pi\eta r v_1 = 0
\end{equation}
Ebből a csepp sugara:
\begin{equation}
r = \sqrt{\frac{3\eta v_1}{2(\sigma-\varrho)g}}
\end{equation}
A q kiszámolásához ezt beírhatnánk az elektron térerősség melletti
mozgásegyenletébe, de ehelyett vonjuk ki a két mozgásegyenletet
egymásból.
Ekkor az $r^3$-os tagok kiesnek:
\begin{equation}
qE = 6\pi\eta r\,(v_1+v_0)
\end{equation}
ide egyszerűbb helyettesítéssel $r$ fenti értékét.

A kondenzátor lemezei közötti térerősséget az $E=U/d$ formulából
számoljuk ki, ahol $U$ a kondenzátorra kapcsolt feszültség, $d$ a lemezek
közötti távolság.
Így kapjuk:
\begin{equation}
q = 9\sqrt{2}\,\pi\,\eta^{3/2}\,(\sigma-\varrho)^{-1/2}\,g^{-1/2}\,
d\,U^{-1}\,v_1^{1/2}(v_1+v_0).
\end{equation}

Ezen összefüggésből a mért esési idők (így az esési sebességek)
alapján Millikan meghatározta az olajcseppek töltéseit.
Eredményül azt kapta, hogy a csepp töltése mindig előállítható egy
egység egész számú többszörösének gondolatával.
Ez az egységet hívjuk elemi töltésnek.
A csepp töltése általában $15$--$20$-szorosa volt az elemi töltésnek,
így az egység hibája nagynak adódott.
Az olaj porlasztásánál nem sikerült kisebb töltésű cseppeket előállítani,
de könnyűnek bizonyult ezt a töltést utóbb egy-két egységgel
megváltoztatni nagy áthatoló képességű radioaktív sugárzás segítségével.
A radioaktív sugárzás ér a levegőt, ionok keletkeznek.
Amikor egy ion a kiszemelt cseppel találkozik, melynek sebességét már
megmértük, az egyensúly hirtelen átalakul és a sebesség megváltozik.

A megváltozott cseppel a fentebb leírt két mérést újra végre kell hajtani.
Az új feszültség $U'$, az új sebességek $u_0$ és $u_1$.
Ekkor a töltésváltozás:
\begin{equation}
\Delta q = K\!\left[ u_1^{1/2}(u_1+u_0){U'}^{-1}
- v_1^{1/2}(v_1+v_0)U^{-1}\right],
\end{equation}
ahol $K$ a fenti egyenletekből adódó állandó:
\begin{equation}
K = 9\sqrt{2}\,\pi\,\eta^{3/2}\,(\sigma-\varrho)^{-1/2}\,g^{-1/2}\,d.
\end{equation}

Ez a képlet csak akkor érvényes, ha egy cseppen sikerül négy mérést
végrehajtani, kettőt besugárzás előtt, kettőt azután.
Millikan megmérte a $\Delta q$ töltésváltozásokat is, és az eredeti
kísérletében kapott elemi érték egy-, két-, háromszorosát mérte a
cseppeken.
Ezzel a töltés egységét — az elemi töltést — pontosabban meghatározta,
és az eredmény egybeesett a Faraday-állandóból az Avogadro-szám
segítségével számolt értékkel, $e$-vel.

Külön probléma, hogy az elektronnak mekkora a töltése. Ezt a következtetést Millikan kísérletéből közvetlenül nem vonhatjuk le, mert ott makroszkopikus cseppek töltését vizsgáltuk.
Az elektronok töltését közvetlenül a sörétzaj vizsgálatával lehetett meghatározni.
Az eredmény, mint ismeretes, az volt, hogy az elektronok töltése
egyenlő az ionokra meghatározott elemi töltéssel.
Az elemi töltés ma elfogadott értéke:

\begin{equation}
e = (1.60217733 \pm 0.00000048)\times 10^{-19}\ \mathrm{C}
\end{equation}

\subsubsection{Sörétzaj}
A sörétzaj, vagy Schottky-zaj, az anódáram először izzókátodos
diódákban megfigyelt ingadozási jelensége.
Az izzókátodos dióda felépítése nagyon hasonlít a katódsugárcsőhöz,
benne a felfűtött katódról az elektronok a termikus energiájukat
felhasználva lépnek ki, és a csövön átrepülve felgyorsulva jutnak
az anódra, így biztosítva az áramvezetést a csövön.
Felerősítés után a hangszóróban pattogó zaj hallatszik, mely a
sörétszemek koppanására emlékeztet.
Oka az, hogy a katódból kilépő elektronos töltést elektronok hordozzák,
és azok véges töltéssel bírnak.
Adott időegység alatt kilépő elektronok száma nem mindig ugyanakkora,
statisztikusan ingadozik.
Ezért az anódáram pontos mérése esetén annak ingadozását is lehet mérni.
Ebből határozta meg W. Schottky az elektron töltését (és így az elemi
töltést), a Millikan-kísérletnél nagyobb pontossággal.

\subsection{Davisson–Germer-kísérlet}
\subsubsection{De Broglie-féle hullámok}
Louis de Broglie 1924-ben vetette fel azt a gondolatot, hogy az anyag ha a fény, ami egy elektromágneses hullám, képes észecske tulajdonságokat mutatni, akkor mivan ha az alapvetően részecskéknek gondolt
anyagok is képesek hullámként viselkedni? Ebből az ötletből vezette le, hogy ha a egy foton energiája $E=h\nu$ és impulzusa $p=E/c=h\nu/c$ akkor a foton hullámhossza:
\begin{equation}
\lambda \nu = c \quad \rightarrow \quad \lambda = \frac{c}{\nu} = \frac{h}{p}
\end{equation}
Ezt az összefüggést pedig alkalmazhatjuk a fénynél lassabban mozgó részecskékre is:
\begin{equation}
\lambda = \frac{h}{p} = \frac{h}{mv}
\end{equation}
Ezt az összefüggést de Broglie hullámhosszának nevezzük. Ez a hullámhossz normális körülmények között mérhetetlenül kicsi. Például egy 1 tonnás autónak, mely 100 km/h-val halad, a de Broglie hullámhossza:
\begin{equation}
\lambda = \frac{6.63 \times 10^{-34} \mathrm{Js}}{(1000 \mathrm{kg})(27.8 \mathrm{m/s})} = 2.38 \times 10^{-37} \mathrm{m}
\end{equation}
Ez a hullámhossz olyan kicsi, hogy semmilyen kísérleti módszerrel nem lehet kimutatni. Azonban a mikroszkópilus világban el lehet érni olyan sebességeket ahol ez a hullámtermészet elvileg
kimutatható kell hogy legyen.

Például ha egy elektront gyorsítunk $U$ feszültséggel, akkor az elektronok $\frac{1}{2}mv^2 = eU$ energiát kapnak (ahol $e$ az elektron töltése), amiből az impulzus:
\begin{equation}
p = mv = \sqrt{2meU}
\end{equation}
Így az elektron de Broglie hullámhossza:
\begin{equation}
\lambda = \frac{h}{\sqrt{2meU}} = \frac{h}{\sqrt{2me}} \cdot \frac{1}{\sqrt{U}} = \frac{1.227}{\sqrt{U}} \; \mathrm{nm}
\end{equation}
Ahol $m = 9.11 \times 10^{-31} \mathrm{kg}$ és $e = 1.602 \times 10^{-19} \mathrm{C}$. Tehát például egy 54 V feszültséggel gyorsított elektron hullámhossza:
\begin{equation}
\lambda = \frac{1.227}{\sqrt{54}} = 0.167 \; \mathrm{nm}
\end{equation}
Ez a hullámhossz már összehasonlítható a kristályok rácsállandójával, így a kristályok rácsai képesek lesznek diffrakciót létrehozni az elektronokon.

\subsubsection{Davisson-Germer kísérlete}
Azt, hogy ez a sejtés akkurrátus-e, először Davisson és Germer mutatta ki 1927-ben, elektronnyalábot vetve egy nikkel egykristályra. Az elektron feltételezett hullámhosszát
a De Broglie sejtés alapján $U$-val szabályozták.
\begin{figure}[H]
    \centering
    \includegraphics[width=0.6\textwidth]{imgs/10-tetel/davisson_germer.png}
    \caption{Davisson-Germer kísérlet elrendezése}
\end{figure}

A feltételezés az, hogy a detektoron egy diffrakciós mintázatot kellene látni, ami igazolná az elektron hullámtermészetét. Ahhoz, hogy a várt diffrakciós képet megértsük,
írjuk fel a szabad elektronhoz rendelhető feltételezett hullámegyenletet. Az elektron az elmélet szerint $\lambda$ hullámhosszal és $\nu = \frac{1}{\tau}$ frekvenciával rendelkezik,
ahol $\tau$ az elektronhoz tartozó periódusidő. Ekkor a hullámegyenlet:
\begin{equation}
\Psi(x,t) = A e^{2i\pi(\frac{x}{\lambda} - \frac{t}{\tau})} = \Psi(x + \lambda, t + \tau) 
\end{equation}
De a hullámot ki lehet fejezni az anyagi részecskére jellemző energiával és impulzussal is, ha felhasználjuk a De Broglie összefüggést ($\lambda = \frac{h}{p}$) és a Planck és Einstein által megadott
energia-frekvencia összefüggést ($E = h\nu = \frac{h}{\tau}$):
\begin{equation}
\Psi(x,t) = A e^{2\pi i (x-Et)/h} = A e^{\frac{i}{\hbar}(px - Et)}
\end{equation}
Ahol $\hbar$ a Dirac által bevezetett jelölés:
\begin{equation}
\hbar = \frac{h}{2\pi} = 1.0545718 \times 10^{-34} \mathrm{Js} = 6.582119569 \times 10^{-16} \mathrm{eVs}
\end{equation}
A koherens forrásból származó elektronok hulláma a szóródás után:
\begin{equation}
\Psi_{s}(r,t) = A e^{\frac{i}{\hbar}(px - Et)} + A e^{\frac{i}{\hbar}(p\left[x + d\right] - Et)}
\end{equation}
Mivel egyes hullámok a kristályrács más-más síkjairól verődnek vissza és mivel $d = d_1 + d_2$ utat tesz meg, így a visszavert hullámok között fáziskülönbség lesz.
\begin{figure}[H]
    \centering
    \includegraphics[width=0.5\textwidth]{imgs/10-tetel/davisson_germer2.png}
    \caption{Koherens elektronnyaláb szóródása vékony fólián. A rétegek távolsága d.}
\end{figure}

Azt már egy ideje tudták, hogy egy kristály rácsra beeső hullám akkor szenved szabályos visszaverődést, ha a Bragg-feltétel teljesül:
\begin{equation}
2d \sin \theta = n \lambda \quad n = 1, 2, 3, ...
\end{equation}
Mivel az intenzitás a hullám amplitúdójának négyzetével arányos, így a visszavert hullám intenzitása:
\begin{equation}
I \propto |\Psi_{s}(r,t)|^2 = 2A^2 \left[1 + \cos\left(\frac{pd}{\hbar}\right)\right]
\end{equation}
A kísérlet során Davisson és Germer a $d$ távolságot nem tudta variálni illetve a beesési szöget is fixen tartották, ezért az egyetlen dolog amit változtatni tudtak az az elektron impulzusa volt a feszültség által.
Ez alapján, és a Bragg-feltétel alapján az intenzitás változásában azt várjuk, hogy:
\begin{equation}
    \cos \frac{pd}{\hbar} = \cos 2\pi \frac{p}{h} d = \cos 2\pi \frac{d}{\lambda} = \begin{cases}
1, & \text{ha } 2\pi \frac{d}{\lambda} = 2\pi n \rightarrow \lambda = \frac{d}{n} \rightarrow I = 4A^2 \\
0, & \text{ha } 2\pi \frac{d}{\lambda} = \frac{\pi}{2} (2n + 1) \rightarrow \lambda = \frac{4d}{2n + 1} \rightarrow I = 2A^2 \\
-1, & \text{ha } 2\pi \frac{d}{\lambda} = \pi (2n + 1) \rightarrow \lambda = \frac{2d}{2n + 1} \rightarrow I = 0 (!)
    \end{cases}
\end{equation}
Tehát a kísérletben azt várjuk, hogy az intenzitás maximumokat akkor kapjuk, amikor a De Broglie hullámhossz kielégíti a Bragg-feltételt. Davisson és Germer a kísérletben
pontosan ezt az eredményt kapta, ezzel igazolva De Broglie sejtését az anyag hullámtermészetéről.

\subsection{Stern–Gerlach-kísérlet}
Az 1920-as évek elejéig a Bohr-féle atommodell volt az, amelyik a legjobban magyarázta a kisérleti eredményeket. Ebben az az atommag egy nagy tömegű, pozitív töltésű pontszerű részecske volt, amely körül az elektronok meghatározott pályákon keringtek.
A Klasszikus elektrodinamika alapján ezeknek a keringő elektronoknak van egy mágneses momentumuk is, melynek a nagysága az áramerősség és a körpálya területének szorzata:
\begin{equation}
\mu = I \cdot A = \frac{e v}{2\pi r} \cdot \pi r^2 = \frac{e}{2m} mvr = \frac{e}{2m} L, \qquad \boldsymbol{\mu} = \frac{e}{2m} \mathbf{L}
\end{equation}
Tehát a keringő elektron mágneses momentumának nagysága arányos a perdületével. A Bohr-féle atommodell szerint az elektronok csak meghatározott értékű perdülettel rendelkezhetnek, így a mágneses momentumuk is kvantált lesz:
\begin{equation}
    \mu = \frac{e}{2m} n\hbar = n \mu_B, \qquad \text{ahol,} \quad \mu_B = \frac{e\hbar}{2m} = 9.27 \times 10^{-24} \; \mathrm{J/T}
\end{equation}
Ez a mennyiség az úgynevezett Bohr-magneton. Mivel ez és a perdület is vektormennyiség, ezért a klasszikus értelmezés szerint a tér bármelyik irányába álhatnak.

\subsubsection{Iránykvantáltság mérése}
Ahhoz, hogy ezt a mágneses momentumot kimutassuk, Stern és Gerlach 1922-ben egy kísérletet tervezett, amelyben több féle atomot egy inhomogén mágneses térbe lőtt be.
A kísérlet lényege, hogy a tér egyik irányát kijelölve (ez tradicionálisan a $z$ irány lesz), mágneses teret hozunk létre egy elektromágnessel
és az ezen átrepülő részecskének az $v_x$ komponense nem változik mivel ebben az irányban nincs erő, azonban a $v_z$ komponensre hat egy erő amit a mágneses momentum és a külső mágneses tér interakciója okoz. 
\begin{figure}[H]
    \centering
    \includegraphics[width=0.6\textwidth]{imgs/10-tetel/stern_gerlach.png}
    \caption{Stern-Gerlach kísérlet elrendezése}
\end{figure}

A mágneses térbe helyezett mágneses momentum helyzeti energiája:
\begin{equation}
    E_{m} = -\boldsymbol{\mu} \cdot \mathbf{B}
\end{equation}
Ebből az erő:
\begin{equation}
    \mathbf{F} = -\nabla E_{m} = \nabla (\boldsymbol{\mu} \cdot \mathbf{B}) = \mu_z \frac{\partial B_z}{\partial z} \hat{z} \quad \rightarrow \quad \mathbf{F} = (0, 0, \mu_z \frac{\partial B_z}{\partial z})
\end{equation}
Ez az erő időben állandó, vagyis egyenletesen gyorsítja a részecskét a $z$ irányban. A szögeltérülés a következőképpen számolható ki:
\begin{figure}[H]
    \centering
    \includegraphics[width=0.4\textwidth]{imgs/10-tetel/stern_gerlach2.png}
    \caption{Szögeltérülés}
\end{figure}
\begin{equation}
    \tan \phi = \frac{v_z}{v_x} \quad \rightarrow \quad \phi = arctan \frac{v_z}{v_x}
\end{equation}
Ahol $v_x$ a részecske kezdeti sebessége a mágneses térbe lépéskor, $v_z$ pedig a mágneses térben szerzett sebesség komponens a $z$ irányban. A kísérlet során Stern és Gerlach
egy speciális ernyőn tette a beérkező részecskéket láthatóvá, és azt nézték, hogy a nulla ponttól mérve (ahová a részecskéknek érkezniük kellene mágneses tér nélkül) mekkora eltérést szenvednek el a részecskék.
Ezt, hogy ha az ernyő $s$ távolságra van és a kitérés $d$ akkor a következőképpen számolhatjuk ki:
\begin{equation}
    \frac{d}{s} = \tan \phi = \frac{v_z}{v_x} = \frac{F_z}{m}\frac{t}{v_x} = \frac{\mu_z \frac{\partial B_z}{\partial z}}{m} \cdot \frac{l}{v_x^2}
\end{equation}
Ahol $l$ a mágneses tér hossza, $m$ a részecske tömege és $t = \frac{l}{v_x}$ a mágneses térben töltött idő. Tehát, ha a $v_x$-et állandónak tartjuk, akkor a kitérés arányos lesz a mágneses momentum $z$ komponensével:
\begin{equation}
    d = \frac{l s}{m v_x^2} \cdot \mu_z \frac{\partial B_z}{\partial z}
\end{equation}
A $\frac{\partial B_z}{\partial z}$-t kísérletileg meg tudták határozni, így a kitérésből a mágneses momentum $z$ komponensét ki tudták számolni.

A kísérletet több részeckével is elvégezték és összefoglalva a következő eredményeket kapták: 
\begin{itemize}
    \item A monoenergiás nyaláb vagy nem tértül el, vagy kevés szám ú komponensre bomlik.
    \item A komponensek távolsága ($\Delta$) egyenlő és szimmetrikus a középvonalra nézve.
    \item Az egyes komponensek intenzitása egyenlő.
\end{itemize}

Az eredményeket úgy lehet értelmezni, hogy a mágneses momentum
nem állhat a tér minden irányában, ahogy azt eddig gondoltuk, hanem
z komponense csak meghatározott értékeket vehet fel az eltérüléssel összahangban.
Ezt hívjuk iránykvantálásnak.
Például a $He$-atomok nyalábja nem bomlott fel, mert a hélium két
elektronjának összesen $0$ eredő mágneses momentuma van.
A $H$ és az $Ag$ két komponensre bomlik, az egyik komponens $\Delta/2$-vel
az el nem térülő pozíció felett, a másik ugyanennyivel alatta helyezkedik el.
Ekkor a mágneses momentum a külső térhez képest két irányban állhat.
A nitrogénatomok nyalábja öt részre bomlott fel, egy komponens középen,
kettő-kettő alul és felül, azonos távolságokra.
Az egyes komponensekben a $\mu_z$ értéke determinált és komponensenként
változik, az eltérülésszerkezet által mutatott kép alapján az egyes
komponensek eltérüléseinek nagysága $\delta_m=m\Delta$ formulával
foglalható össze, ahol $m$ egy egész vagy félegész szám (pozitív vagy
negatív annak megfelelően, hogy felfelé vagy lefelé történt az eltérülés),
$\Delta$ a komponensek távolsága.
Az eltérülés arányos a $\mu_z$-vel, ezért ugyanez a mágneses
momentum $z$ komponenséről is elmondható.
A $\mu_z$ értékei csak diszkrét értékeket vehetnek fel a következő módon:
\[
\mu_z = m\cdot\mu_0,
\]
ahol $\mu_0$ a $\Delta$ távolság és egy az $\alpha$-eltérülési szög és a
mágneses momentum viszonyát megadó (6.2) egyenletből kiszámolható
konstans szorzata.
Az eltérülésszerkezetek mindig szimmetrikusan, egymástól egyenlő
távolságra helyezkedtek el az eltérülés nélküli foltra nézve, ezért azt is
állíthatjuk, hogy az $\mu_z$ értékei az $m_{\min}$ és az $m_{\max}$
és az $-m_{\max}$ között egész lépésekben változhatnak csak.
Például egy 4 komponensre bomló sugárnál az egyes komponensekhez tartozó $m$ értéke alulról felfelé a következő
módon változik: $-\tfrac{3}{2}, -\tfrac{1}{2}, \tfrac{1}{2}, \tfrac{3}{2}$,
az eltérülés kvantitatív értéke ezek $\Delta$-szorosa.

A mágneses momentum „nagysága” ($\mu$) megállapodás szerint a legnagyobb
$\mu_z$ érték, azaz a legjobban eltérült komponenshez tartozó $\mu_z$.
Ebből és a felhasadt komponensek számából a többi $\mu_z$ is kiszámítható,
amelyek nem a legszélső irányba térülnek el.
Az eltérések kvantitatív vizsgálatából a Stern–Gerlach-kísérletben az
derült ki, hogy hidrogénatom-nyaláb és ezüstatom-nyaláb esetén:

\begin{equation}
\mu_z^{\max} = \mu = -\frac{e\hbar}{2m} = \mu_B.
\end{equation}

A hidrogén esetén csak egy elektron van az atomi héjakon, ezért ez az
elektron felelős az atom mágneses momentumáért.
(Az atommag mágneses momentuma nagyságrendekkel kisebb.)
A Bohr-modellben, az $n=1$-es főkvantumszámú pályán kört pályán keringő
(alapállapotú) elektron klasszikusan értelmezett mágneses momentuma
éppen ennyinek adódott.
Azt viszont már nem tudjuk a Bohr-modell keretein belül értelmezni,
hogy miért két részre hasad fel a hidrogénatom-nyaláb.

\subsubsection{Több mágneses tér használata}
\begin{figure}[H]
    \centering
    \includegraphics[width=0.6\textwidth]{imgs/10-tetel/stern_gerlach3.png}
    \caption{Több Stern-Gerlach mágnes egymás után}
\end{figure}
Egy kiválasztott komponens ugyanolyan irányú újraanalizálásakor már nem
bomlik fel még egyszer.
Ha azonban a második mágnes $90^\circ$-kal el van forgatva úgy, hogy az
inhomogenitás y irányú, és megfelelően gyorsan továbbítjuk a nyalábot
ezen második mágnesbe, akkor a kiválasztott komponens tovább bomlik
nem azonos intenzitású, de az előzővel megegyező eltérülésszerkezetre.
Ebből azt a következtetést vonhatjuk le, hogy a $\mu_z$ nem lehet
akkora, mint a mágneses momentumvektor hossza, ugyanis a legnagyobb z
irányú eltérüléshez tartozó z komponens is fel tudott még bomlani y
irányban, így van $\mu_y$ komponense is.
Ezért a mágneses momentumvektor abszolút értéke nagyobb, mint $\mu_z$!
A Bohr-modell jól adta meg a mágneses momentum harmadik komponensének
maximális értékét, de ez nem a mágneses momentumvektor nagysága!
A Bohr-modell érvényességi köre itt már véget ér.

Ha három Stern–Gerlach-mágnest alkalmazunk egymás után, akkor érdekes
jelenséget tapasztalunk.
Válasszuk ki az első mágnessel a z irányú egyik komponenst.
Ezután analizáljuk ezt egy y irányba elforgatott mágnessel, ilyenkor,
mint említettük, ismét komponensekre bomlik az adott $\mu_z$ értékű nyaláb.
Ha ezután ismét egy z irányú mágnessel analizáljuk a már két mágneses
téren áthaladt nyalábot, akkor azt várnánk, hogy az csak egy komponenst
mutat, hiszen a z szerint már megtörtént az egy komponens kiválasztása.
A kísérlet azt mutatja, hogy ilyenkor a harmadik mágnes után is több
csoportot lehet tapasztalni!
Ez a kvantummechanikai méréselmélet fogalomkörében magyarázható.
A mérés beavatkozás a rendszerbe, néha egy mennyiség mérése más mennyiségek
értékeit megváltoztatja.

Fontos megjegyeznünk, hogy a Bohr-féle kvantumfeltétel a perdület, és így
a mágneses momentum nagyságát kvantálja, az alapállapotú atomokra egy jól
meghatározott perdületet és mágneses momentumot eredményez.
Ennek megfelelően a Stern–Gerlach-kísérletben a mágneses momentum abszolút
értéke, és így annak négyzete, minden komponensben megegyezik.
Az újdonság itt egy külső tér által kijelölt irányra vett vetület értékeinek
kvantumos viselkedése, ezt neveztük iránykvantálásnak.
Megállapodás szerint a külső tér által kijelölt irányban vesszük fel a
koordináta-rendszer z tengelyét.
(Gyakran egy vektor z komponensét a vektor harmadik komponensének is
nevezzük.)

A Stern–Gerlach-kísérlet alkalmas molekulák eltérítésére is.
Zárt elektronszerkezetű molekuláknál az eltérülést az atommag mágneses
momentuma okozza.
Ezzel a módszerrel mérték meg először a proton mágneses momentumát,
ami az elmélettől (Dirac-egyenlet) eltérő értéknek adódott.



\subsection{Einstein–de Haas-kísérlet}
A Stern–Gerlach-kísérlet a mágneses momentum külső tér irányú komponensét mérte meg.
Mint láttuk, a mágneses momentum összefüggésben van a perdülettel.
Felmerülhet a kérdés, hogy a perdület és a mágneses momentum kapcsolatát
klasszikusan megadó arányosságot a mérések igazolják-e.

Az Einstein–de Haas-kísérlet (1915) az elektronok mágneses momentumának
átfordításával járó perdületváltozást méri.
A kísérletben egy ferromágneses hengert áramjárta tekercsel felmágneseznek,
az áram irányát hirtelen megváltoztatják.
Ilyenkor a mágnesben az elektronok, mint kis iránytűk, először beállnak
párhuzamosan, jó nagy eredő mágneses momentummot kialakítva, majd az
áramirány-változásra az összes átfordul, magával fordítva a perdületeket is.
A perdületváltozás/Idő azonban forgatónyomatékot jelent, ezért a vékony
torziós szálon függő henger elfordul.
A szálon egy tükröt vékony fénynyalábbal megvilágítva a tükörkép helyzete
jelzi, hogy a henger mennyit fordult el.
Az átmágnesezést periodikusan a torziós ingának lengésidejének megfelelően
beállítva, rezonanciaszerűen végezték.
\begin{figure}[H]
    \centering
    \includegraphics[width=0.6\textwidth]{imgs/10-tetel/einstein_de_haas.png}
    \caption{Einstein–de Haas-kísérlet elrendezése}
\end{figure}
A kísérlet eredménye, ahogy E. Back gondos mérései is mutatták (1919),
hogy a mágneses momentum és a perdület aránya
\[
\frac{e}{m},
\]
a várt érték helyett
\[
\frac{e}{2m}.
\]
A \(\dfrac{e}{m}\) egy körpályán mozgó töltésre volt igaz.
Az a kísérleti tény, hogy az elektronok az Einstein–de Haas-kísérlet
ferromágnesében kétszer akkora mágneses momentumot mutattak, mint a
körpálya miatt kellene, sokáig megoldatlan probléma maradt.

A giromágneses faktor az arányossági tényező a mikrorészecskék perdülete
($\vec{I}$) és mágneses momentuma között:
\begin{equation}
\vec{\mu} = \gamma \vec{I}.
\end{equation}

A keringő elektron esetén ez $\gamma=\dfrac{e}{2m}$, de ebben a
kísérletben $\gamma=\dfrac{e}{m}$, amely a helyes mágneses momentumot adja.
A giromágneses faktornak van dimenziózása, ezért vezessünk be egy
dimenzió nélküli mennyiséget, ami lényegében megfelel a giromágneses faktornak.
Ez a g-faktor, vagy dimenziótlan giromágneses faktor, jele $g$.
A g-faktor használatával a fenti arányosság a következő módon írható:

\begin{equation}
\vec{\mu} = g\,\frac{e}{2m}\,\vec{I}.
\tag{6.3}
\end{equation}

Szemléletesen ez azt jelenti, hogy ha a mágneses momentumot az ő
természetes egységében, $\mu_B=\dfrac{e\hbar}{2m}$-ban, a perdületet
pedig $\hbar$-ban mérjük, akkor az arányossági tényező a g-faktor:
\begin{equation}
\frac{\vec{\mu}}{\mu_B} = g\,\frac{\vec{I}}{\hbar}.
\end{equation}

1921-ben A. Compton vetette fel, és 1925-ben Goudsmit és Uhlenbeck
mutattak rá (az eddig elmondott kísérleteken kívül még az anomális
Zeeman-effektus tapasztalatai alapján), hogy az elektronoknak saját
forgásából adódó mágneses momentuma is lehet, méghozzá kétszeres
giromágneses faktorral.
A kvantummechanika relativisztikus egyenletének, a Dirac-egyenletnek a
felállítása után pedig egyszerűen az egyenletből adódik az elektron ezen
tulajdonsága, hogy van egy perdülete, és ebből adódóan mágneses momentuma
mindenféle mozgás nélkül.
Ezt sajátperdületnek hívjuk, vagy más néven spinnel.
A Dirac-egyenlet megoldásából éppen az $e/m$ arány adódik a mágneses
momentum és a sajátperdület arányára, avagy a g-faktorra, 2.
Így az Einstein–de Haas-kísérletben a ferromágnességet az elektronok
saját mágneses momentumainak tulajdonítjuk.
Ez a kísérlet az elektronspin kétszeres mágnesességének első egyértelmű
bizonyítéka.

Felmerül a kérdés, hogy az alapállapotú hidrogénatom elektronja milyen
perdülettel rendelkezik.
A két lehetséges válasz: a klasszikusan is ismert pályaperdület, vagy az
Einstein–de Haas-kísérletben is tapasztalt sajátperdület.
Ha nincs sajátperdülete, csak pályaperdülete, akkor a pályaperdület
$1 \cdot \hbar$ értéke mellett (Bohr-modell) megkapjuk a Stern–Gerlach-kísérlet
eredményét, a $\mu=\mu_B\cdot 1$-et.
Ha csak sajátperdülete van, akkor kétszeres mágnesessége is van hozzá,
és $\mu=2\mu_B\cdot\frac{1}{2}$ szerint kapjuk vissza a Stem- Gerlach-kísérlet eredményét, tehát ilyenkor a perdületre a Bohr-modellel ellentétben csak $\hbar/2$ marad. A két lehetséges válasz között egy
újabb mérés dönt. A pereltilet nagysága a kvantummechanika szerint összefügg azzal, hogy hány részre bomlik fel az atomnyaláb mágneses térben
(Stern-Gerlach-kísérlet) vagy egy színképvonal mágneses tér ben (Zeeman-effektus). A felhasadások száma dönti el, hogy a hidrogénatom elek tronjának milyen fajta perdülete van.

\subsection{Zeeman-effektus}
A XIX. század végén nagyszámú kísérleti tapasztalat halmozódott fel az
atomok színképeiről.
A színképek nagyon hasznos eszköznek bizonyultak, mert az elemek
meghatározását tették lehetővé.
A vonalak szerkezetének értelmezése nehéz feladatnak bizonyult, és a
klasszikus fizikával nem magyarázható kvantumhipotézis kellett hozzá
(Bohr-féle kvantumfeltétel).
A színképvonalak mágneses térben történő viselkedése is alapvetően új
gondolatokat hozott.
Ezt kísérletben először P. Zeeman vizsgálta 1897-ben.
Ő megmutatta, hogy a színképvonalak mágneses térben felhasadnak: ez a
Zeeman-effektus.
A felhasadás sokszor nagyon bonyolult színképstruktúrákat szolgáltat,
de az esetek egy részét sikerült a klasszikus elektronelmélettel az
atomok Thomson-modelljében megmagyarázni.
A klasszikusan is magyarázható szerkezetet normális Zeeman-effektusnak
hívjuk.
A normális felhasadást Lorentz előre meg is jósolta 1896-ban
(1902-ben közösen kaptak érte Nobel-díjat).
A többi típusú felhasadást anomális Zeeman-effektusnak hívjuk, ennek
magyarázatához a spin fogalma alapvetően hozzátartozik.

A spektroszkópia technikai fejlődésével az egyes vonalak finomabb
szerkezete is megfigyelhetővé vált: Voltak ún. szingulett vonalak,
melyeknek mágneses tér nélkül nem volt finomszerkezete.
Más vonalak két, egymáshoz nagyon közel álló vonal együtteséi voltak,
ezeket hívjuk dublettnek, értelmezésükről a következő fejezetekben lesz szó.
A színképvonalak ezekbe nem sorolható részét hívjuk multipletnek.
A normális Zeeman-effektust a szingulet vonalak mutatják, anomális
Zeeman-effektust a dublett és multiplet vonalaknál tapasztaltak.

A Zeeman-effektus mérésére szolgáló kísérletben a gerjesztett atomokat (legegyszerűbben gázt) két mágneses
E–D pólusai között kialakított, jó közelítéssel homogén mágneses térbe
helyezzük.
A $\mathbf{B}$ mágneses tér irányát hívjuk z iránynak.
Ilyenkor a z irányra merőlegesen kilépő fénysugarakat egy
kollimátorral kiválasztjuk, és prizmára vagy optikai rácsra (R) ejtve a
színkép tanulmányozható.
Normális esetben a mágneses tér nélkül $\nu_0$ frekvenciánál található
szingulett vonal felhasad három részre.
A $\nu_0$ mellé nála $\pm \Delta \nu$ frekvenciával magasabb és alcsonyabb frekvenciájú vonalak kerülnek (Lorentz-féle triplett). Mindhárom
komponens síkban polarizált fénysugár. Az eredeti frekvenciánál tapasztalt
vonal a mágneses tér irányában polarizált ún. $\pi$-komponens, a másik kettő
az x és y irányban polarizált $\sigma$-komponensek. A tapasztalat és az elmélet
szerint is a felhasadáskor a $\Delta \nu$ eltolódás arányos a mágneses tér nagyságával.
\begin{figure}[H]
    \centering
    \includegraphics[width=0.6\textwidth]{imgs/10-tetel/zeeman.png}
    \caption{A Zeeman-féle kísérlet sematikus elrendezése és eredményeinek szemléltetése. Az a) esetben mágneses tér nélkül $\nu_0$ frekvenciájú fényt detektálunk. A b)
esetben a mágneses térre merőlegesen megfigyelve három vonaira hasad a színkép. A c) esetben a mágneses térrel párhuzamosan kilépő fénysugarak csak két frekvenciájú részre bomlanak (ezek cirkulárisan polarizált komponensek lesznek)}
\end{figure}
A normális Zeeman-effektus klasszikus magyarázatának alapja az, hogy az
atom \(\vec{\mu}\) mágneses momentumára homogén mágneses térben
\(\vec{M}\) forgatónyomaték hat, emiatt az atom egész elektronrendszere
forgásba jön, így a mágneses momentum is.
A forgásra felírt perdülettétel most így alakul:
\begin{equation}
\frac{1}{\gamma} \frac{d\vec{\mu}}{dt} = \frac{d\vec{I}}{dt}=\vec{M}=\vec{\mu}\times\vec{B}.
\end{equation}
Vagy átalakítva:
\begin{equation}
\frac{d\vec{\mu}}{dt}=\gamma\,\vec{\mu}\times\vec{B}.
\end{equation}
Ennek megoldása az, hogy a \(\vec{\mu}\) az atom mágneses momentuma és
a vele párhuzamos \(\vec{I}\) perdület a \(\vec{B}\) mágneses indukció
iránya körül
\begin{equation}
\omega_L=\gamma B
\end{equation}
szögsebességgel precessziót végez, azaz az atomban nyugvó elektronok
\(x,y\) síkban forognak ezzel a frekvenciával.
A
\begin{equation}
\nu_L=\omega_L/2\pi
\end{equation}
frekvenciát Larmor-frekvenciának nevezzük.
Amikor az atom fényt bocsát ki, azt most úgy értelmezzük, hogy az atomi
elektron \(\nu\) frekvenciájú oszcillációt végez, \(\vec r'\) helyvektora
időben \(\vec r'=\vec a\cos\omega t\) szerint rezeg, ezért sugároz.

Ha belül az elektron \(\nu\) frekvenciával rezeg, valamint az atom
\(\omega_L\) szögsebességgel forog, akkor a rezgés és a forgás
\(x\) és \(y\) irányú harmonikus rezgésekre bonthatók.
Mindkét irányban a két rezgés összegzésekor egy \(\nu+\Delta\nu\) és
egy \(\nu-\Delta\nu\) frekvenciájú rezgés összegét kapjuk.
Végezzük ezt el az \(x\) irányra.
Az atomon belül \(\mathbf{r}'\) helyvektor mutat az elektron helyére.
Ennek \(x\) komponense a térbeli polárkoordináták felhasználásával
\(x = r'\sin\vartheta\cos\varphi\).
A forgás és a rezgés miatt
\(r' = a\cos\omega t\) és \(\varphi=\omega_L t\), ezért trigonometrikus
azonosságokkal átalakítva kapjuk:

\begin{equation}
x = \frac{a}{2}\sin\vartheta\big(\cos(\omega+\omega_L)t + \cos(\omega-\omega_L)t\big).
\end{equation}

Az \(y\) irányból észlelve a fényt \(x\) és \(z\) irányban polarizált
módusokra bonthatjuk. (Részletesen lásd a következő alfejezetben.)
Az \(x\) irányban polarizált módus (\(\sigma\)) azonban az atom
precessziója miatt két eltolt frekvenciájú rezgés összege lesz.
A \(z\) irányban polarizált \(\pi\)-módus eltolás nélkül észlelhető,
azt a mágneses tér nem módosítja.
Ez a modell a pályaperdületből számolt mágneses momentumokkal dolgozik
(1896-ban Lorentz dolgozta ki), és kizárólag akkor használható, ha az
atomban az elektronok sajátperdületei kompenzálják egymást,
csak pályaperdületből származó mágneses momentum van
(ezek a szingulett vonalak).

Amikor magános elektron okozza az atom mágneses momentumát, akkor
anomális felhasadást tapasztalunk.
Erre példa a nátrium D-vonal-dublettjének kisebb frekvenciájú tagja,
ami négy vonalra hasad, két vonalpárosra.
Az anomális Zeeman-felhasadást a klasszikus fizika nem tudta megmagyarázni.
Ez volt az egyik alapvető tapasztalat, amely az elektron sajátperdületének,
vagy spinjének gondolatához vezetett.
A Zeeman-effektus teljes magyarázatához a kvantummechanika segítségével
lehetőségessé vált.

Ezt a vonalfelhasadást akkor nevezzük Zeeman-effektusnak, ha a külső
mágneses tér nagysága kisebb az atomon belüli mágneses tereknél.
Ilyenkor a Zeeman-felhasadás a multiplet vonalak távolságánál kisebb marad.
Ha a külső mágneses tér miatti vonalfelhasadás a multipletek vonalainak
távolságát is eléri vagy meghaladja, akkor nagy mágneses térről beszélünk.

Az ilyen nagy külső tér esetén a vonalak felhasadásának szerkezete
megváltozik, nagy mágneses térben a pálya- és a sajátperdületekből
származó mágneses momentumok belső terekből adódó szoros kapcsolata
felbomlik.
A Zeeman-effektusnál a pálya- és a sajátperdület mint egység vett fel
különböző irányítottságú állapotokat.
Nagy mágneses térben azonban a kétféle mágneses momentum közötti kapcsolat
energianyere­ségét — amit az atomon belüli mágneses terek okoznak — legyőzi
a külső tér hatása, energetikailag nyereséges lesz a spinnek a pályaperdülettől
elfordulnia (szétcsatolódnak).
A gerjesztett atom ilyen esetben úgy is adhat le egy fotont, hogy csak a
pályaperdületből származó mágneses momentuma változik, vagy forditva.
A pályaperdület z komponense és a sajátperdület z komponense egymástól
függetlenül vehetik fel értékeiket.
Így az anomális felhasadásból is normális felhasadás lesz.
Ez a Paschen–Back-effektus (1912).

\begin{figure}[H]
    \centering
    \includegraphics[width=0.6\textwidth]{imgs/10-tetel/zeeman2.png}
    \caption{A Zeeman effektus szemléletetése}
\end{figure}

\end{document}